% ---------------------------------------------------------------------
% --- Arquivo principal e os demais serao os dos capitulos.
% --- EXPRESSÔES ENTRE <> DEVERÂO SER COMPLETADAS COM A INFORMAÇÂO ESPECÌFICA DO TRABALHO 
% ---------------------------------------------------------------------
\documentclass[ruledheader]{abnt_UFF}[abntex2]
\usepackage[alf]{abntex2cite}

%--- pacotes para hiphenizacao e acentuacao em portugues
\usepackage[brazil]{babel}
%\usepackage[latin1]{inputenc}
\usepackage[utf8]{inputenc}
\usepackage[T1]{fontenc}

%--- pacote para figuras
\usepackage{epsf}
\usepackage[dvips]{epsfig,graphicx}
\usepackage{subfigure}
\graphicspath{{imagens/}}

%--- pacote de simbolos
\usepackage{latexsym}
\usepackage{textcomp}

%--- simbolos matematicos
\usepackage{amsmath}
\usepackage{amssymb}

%--- pacote para gerar pseudo-codigo
\usepackage{algorithm}
\usepackage{algorithmic}
\floatname{algorithm}{Algoritmo}

%--- outros pacotes
\usepackage{url}
\usepackage{longtable}
\usepackage{pdflscape}
\usepackage{setspace}
\usepackage{booktabs}
\usepackage{alltt}
\usepackage[flushleft]{threeparttable}
\renewcommand{\ttdefault}{txtt}

%Tabela Colorida
\usepackage{colortbl}
\usepackage{multicol}
\usepackage{multirow}
\usepackage{rotating}

% Listas personalizadas
\usepackage{float}
\newfloat{quadro}{htbp}{loq}[chapter]
\floatname{quadro}{Quadro}
\newcommand{\listofquadros}{\listof{quadro}{Lista de Quadros}}

\newfloat{grafico}{htbp}{logr}[chapter]
\floatname{grafico}{Gráfico}
\newcommand{\listofgraficos}{\listof{grafico}{Lista de Gráficos}}

\hyphenation{
    a-de-qua-da-men-te
    di-men-sio-na-men-to
}

%---------usando tipo de fonte padrao
\renewcommand{\ABNTchapterfont}{\bfseries\fontfamily{cmr}\fontseries{b}\selectfont}
\renewcommand{\ABNTsectionfont}{\bfseries\fontfamily{cmr}}

% --- -----------------------------------------------------------------
% --- Documento Principal.
% --- -----------------------------------------------------------------
\usepackage{nameref}
% \usepackage[pdftex]{hyperref}
% \hypersetup{colorlinks, citecolor=black, pdftex}

\newcommand{\meunome}{CHRISTOPHER RIC COELHO SAAR}
\newcommand{\meutitulo}{EXPLORAÇÃO DE ESTADOS IMPLÍCITOS EM REQUISITOS DE SOFTWARE COMO ESTRATÉGIA PEDAGÓGICA NO ENSINO DE ENGENHARIA DE REQUISITOS E TESTES}
\newcommand{\meuorientador}{Prof. Dr. Daricélio Moreira Soares}
\newcommand{\meuano}{2026}
\newcommand{\meumes}{<MES>}
\newcommand{\meudia}{<DIA>}

\begin{document}
    % --- Para elementos em português, remover.
    \renewcommand{\listfigurename}{Lista de Figuras}
    \renewcommand{\figurename}{Figura}
    \renewcommand{\listtablename}{Lista de Tabelas}
    \renewcommand{\tablename}{Tabela}
    \renewcommand{\contentsname}{Sumário}
    \renewcommand{\refname}{Referências}
    \renewcommand{\appendixname}{Apêndices}
    \renewcommand{\chaptername}{Capítulo}

% --- -----------------------------------------------------------------
% --- Titulo, abstract, dedicatorias e agradecimentos.
% --- Indice geral, lista de figuras e tabelas.
% --- -----------------------------------------------------------------
	% --- -----------------------------------------------------------------
	% --- Elementos usados na Capa e na Folha de Rosto.
	% --- EXPRESS�ES ENTRE <> DEVER�O SER COMPLETADAS COM A INFORMA��O ESPEC�FICA DO TRABALHO
	% --- E OS S�MBOLOS <> DEVEM SER RETIRADOS 
	% --- -----------------------------------------------------------------
	\cleardoublepage
	\thispagestyle{empty}
	
	\vspace{-60mm}
	  \begin{singlespace}
	    \begin{center}
	      {\large UNIVERSIDADE FEDERAL DO ACRE}
	      \vskip4.0cm{\textbf{\Large \meunome}}
	      \vskip5.0cm {\textbf{\meutitulo}}
	    \end{center}
	    \begin{center}
	      \vskip10.0cm{\textbf{RIO BRANCO\\\meuano}}
	    \end{center}
	  \end{singlespace}
	% --- -----------------------------------------------------------------
	% --- Folha de rosto. (Obrigatorio)
	% --- ----------------------------------------------------------------
	\cleardoublepage
	\thispagestyle{empty}
	
	\vspace{-60mm}
	  \begin{center}
	    {\large UNIVERSIDADE FEDERAL DO ACRE}\\
	    \vspace{3cm}
	    {\large \meunome} \\
	    \vspace{3cm}
	    
	    {\large \meutitulo} \\
	    \vspace{1.5cm}
	  \end{center}
	
	\noindent
	  \begin{flushright}
	    \begin{minipage}[t]{8cm}
	      Proposta de dissertação de mestrado submetida ao Programa de Pós-Graduação em Ciência da Computação na Universidade Federal do Acre como requisito parcial para obtenção do título de mestre em Ciência da Computação. Linha de Pesquisa: Engenharia de Software.
	    \end{minipage}
	  \end{flushright}
	
	\vskip1.50cm
	  \begin{center}
	    \small Orientador: \\
	    \meuorientador\vskip3.5cm 
	  \end{center}
	  \begin{center}
	    \vspace{4mm}
	    RIO BRANCO \\
	    %\vspace{6mm}
	    \meuano
	  \end{center}
	
	\thispagestyle{empty}
	% --- -----------------------------------------------------------------
	% --- Termo de aprovacao. (Obrigatorio)
	% --- ----------------------------------------------------------------
	\cleardoublepage
	\thispagestyle{empty}
	
	\vspace{-60mm}
	
	\begin{center}
	  {\large \meunome} \\
	  \vspace{7mm}
	
	  {\large \meutitulo} \\
	  \vspace{10mm}
	\end{center}
	
	\noindent
	\begin{flushright}
	  \begin{minipage}[t]{8cm}
	
		Proposta de dissertação de mestrado submetida ao Programa de Pós-Graduação em Ciência da Computação na Universidade Federal do Acre como requisito parcial para obtenção do título de mestre em Ciência da Computação. Linha de Pesquisa: Engenharia de Software.
	
	  \end{minipage}
	\end{flushright}
	
	\vspace{1.0 cm}
	\noindent
	
	Aprovado em \meudia\ de \meumes\ de \meuano. \\
	\begin{flushright}
	  \parbox{11cm}
	  {
	    \begin{center}
	      \vspace{6mm}
	      \rule{11cm}{.1mm} \\
	      \meuorientador\\
	      Universidade Federal do Acre
	      \vspace{6mm}
	    \end{center}
	  }
	\end{flushright}
	
	\begin{center}
	  \vspace{4mm}
	  RIO BRANCO \\
	  %\vspace{6mm}
	  \meuano
	\end{center}
	
	% --- -----------------------------------------------------------------
	% --- Dedicatoria.(Opcional)
	% --- ----------------------------------------------------------------
	\cleardoublepage
	\thispagestyle{empty}
	\vspace*{200mm}
	
	\begin{flushright}
	  {\em
	    Tudo o que temos de decidir é o que fazer com o tempo que nos é dado.
	    \\
	    Gandalf.
	    \\ 
	    O Senhor dos Anéis - A Sociedade do Anel (2001)
	  }
	\end{flushright}
	\newpage
	
	% --- -----------------------------------------------------------------
	% --- Agradecimentos.(Opcional)
	% --- ----------------------------------------------------------------
	\pretextualchapter{Agradecimentos}
	\hspace{5mm}
	
	A mim e aos meus.
	
	% --- -----------------------------------------------------------------
	% --- Resumo em portugues.(Obrigatorio)
	% --- ----------------------------------------------------------------
	\begin{resumo}
	
	  \begin{center}{
	    \textbf{\meutitulo}}
	  \end{center}
	
	\begin{spacing}{1.0}\noindent

        resumo
    
	\end{spacing}
	{\hspace{-8mm} \bf{Palavras-chave}}: Engenharia de Software.
	
	\end{resumo}
	
	% --- -----------------------------------------------------------------
	% --- Resumo em lingua estrangeira.(Obrigatorio)
	% --- ----------------------------------------------------------------
	\begin{abstract}\noindent 
        abstract
	
	{\hspace{-8mm} \bf{Keywords}}: Software Engineering
	
	\end{abstract}
	
	% --- -----------------------------------------------------------------
	% --- Lista de figuras.(Opcional)
	% --- -----------------------------------------------------------------
	%\cleardoublepage
	\listoffigures
	
	% --- -----------------------------------------------------------------
	% --- Lista de tabelas.(Opcional)
	% --- -----------------------------------------------------------------
	\cleardoublepage
	%\label{pag:last_page_introduction}
	\listoftables
	\cleardoublepage

	% --- -----------------------------------------------------------------
	% --- Lista de quadros.(Opcional)
	% --- -----------------------------------------------------------------
	\cleardoublepage
	\listofquadros
	\cleardoublepage

	% --- -----------------------------------------------------------------
	% --- Lista de gráficos.(Opcional)
	% --- -----------------------------------------------------------------
	\cleardoublepage
	\listofgraficos
	\cleardoublepage
	
	% --- -----------------------------------------------------------------
	% --- Sumario.(Obrigatorio)
	% --- -----------------------------------------------------------------
	\pagestyle{ruledheader}
	\tableofcontents
	
% --- -----------------------------------------------------------------
% --- Insercao dos capitulos.
% --- -----------------------------------------------------------------
    \pagestyle{ruledheader}
    \setcounter{page}{1}
    \pagenumbering{arabic}
\chapter{Introdução}
\label{cap1title}

É comum que pesquisadores busquem entender os efeitos que uma variável pode ou não exercer sobre outras, ainda que não exista entre elas uma relação direta que, no entanto, poderia ser genericamente representada por $y=f(x_{1},x_{2},...,x_{n})$. Desse modo, por exemplo, com base nos valores conhecidos para $x$ é possível estimar a influência dele sobre $y$. Esta ideia representa o conceito dos modelos de regressão que, em outros termos, buscam explicar a relação entre variável dependente (resposta) em função de uma ou mais variáveis explicativas (independentes) \cite{chein2019introduccao, angrist2009mostly}

Em alguns casos, estas relações podem apresentar um esquema estrutural em que a disposição dos dados exigem uma relação de subordinação entre variáveis \cite{natis2001modelos}. Em grande parte dos estudos com essa característica a variável "tempo" é o fator de associação mais comum, de modo que no modelo hierárquico sua relação pode ser descrita genericamente por  $\mathbf{y}_{i}=(y_{i1},\cdots,y_{im_{i}})$, sendo $y_{ik}$ a resposta para o $i$-ésimo indivíduo no $k$-ésimo tempo, $k=(1,\cdots,m_{i})$ \cite{singer2000analysis}.
 

 

A aplicação prática desse conceito pode ser encontrada, por exemplo, ao se perceber a variação de melhora ou piora em um paciente quando submetido a um tratamento ao longo do tempo, ou ainda, quando uma região de solo é submetido a controles que influenciam nos níveis de seus nutrientes para diferentes observações no tempo. Em ambos os casos, o delineamento do modelo deve considerar as variáveis dependentes para cada tempo e assim concluir sobre a eficácia dos controles aplicados.


Estudos em que as observações sobre os indivíduos são coletadas repetidamente ao logo do tempo são classificados pela literatura como \textbf{medida repetida no tempo}, \textbf{dados longitudinais}, \textbf{coorte}, \textbf{painel}, entre outros \cite{singer2000analysis, rosa2001analise, fagundes2013metodologias}.  


%Na análise de regressão é comum se deparar com situações práticas onde existe o interesse de entender o comportamento de uma ou mais variáveis respostas medidas ao longo de uma dimensão ordenada \cite{singer2000analysis}. Um exemplo disso corre na modelagem da variação do carbono no solo com ou sem cobertura vegetal durante um certo ciclo de observações, como foi o caso do pesquisa proposta por \cite{gabriel2022}. 

% Estudos com esse pressuposto, onde há observações repetidas sobre a mesma unidade amostral ao longo do tempo, são categorizados como estudos de medida repetida no tempo, podendo ainda assumir outras nomenclaturas na literatura tais como: longitudinal, coorte, painel, entre outras.  . 


Ainda que sejam comuns, os estudos de medida repetida no tempo apresentam desafios maiores quando comparado a outros ajustes \cite{fitzmaurice2008longitudinal}, resultado da natureza de correlação estrutural dos dados, implicando na necessidade de uma modelagem que satisfaça a condição de subordinação dos grupos \cite{cnaan1997using, fausto2008modelo}. 

% Apesar de comum, pesquisas dessa natureza ainda apresentam desafios, sobretudo quanto a modelagem adequada dos dados . Devido à sua natureza, os valores resultantes das observações repetidas ao logo do tempo apresentam uma estrutura de correlação onde a modelagem desempenha um papel fundamental na análise correta dos dados. \cite{cnaan1997using, fausto2008modelo}. 

Atualmente, o padrão de delineamento de dados longitudinais é dado pelo uso de modelos lineares misto ou de efeito aleatório, \cite{laird1982random}, em que exigi-se a pressuposição de um comportamento gaussiano da variável resposta ou dos erros, de modo que na ausência de tal pressuposto cabe ao pesquisador aplicar uma  transformação nos dados a fim de normalizá-los \cite{cnaan1997using}.

%A abordagem até então predominante no delineamento de dados longitudinais, era a utilização de modelos lineares misto ou de efeito aleatório \cite{laird1982random}, no entanto, este modelos exigiam a pressuposição de um comportamento gaussiano da variável resposta, o que nem sempre era possível, e nestes casos o pesquisador aplicava uma transformação a fim de normalizar o dados \cite{cnaan1997using}.  

Em síntese, os modelos mistos, na ausência de normalidade dos erros, apresentavam uma flexibilidade menor na captura de varições entre a variável resposta, exigindo técnicas de normalização dos dados. Em alguns caso o modelo pode até gerar resultados razoáveis, no entanto, deve ser evitado em detrimento a abordagens mais eficientes \cite{azzalini1999statistical, braga2022random}.

% Um outro fator de impacto direto no tratamento dos dados refere-se a tipo de variável resposta, podendo ela ser contínua ou discreta. Nesse contexto, modelos lineares de efeitos mistos, ou simplesmente modelos de efeito aleatório \cite{laird1982random} tem sido a abordagem predominantemente usadas no tratamento de dados longitudinais quando a variável resposta tem distribuição normal o que, no entanto, em diversos casos, acaba limitando-a a valores de zero e um \cite{cnaan1997using}. Ou seja, na ausência de normalidade dos erros ou efeitos aleatórios, os modelos mistos tendem a apresentar menor flexibilidade para capturar as dissonâncias entre os indivíduos, sendo necessário a transformação da variável resposta, o que em muitos casos leva a resultados até razoáveis, mas deve ser evitado em detrimento de um método melhor \cite{azzalini1999statistical, braga2022random}. 

Os aspectos descrito acima, demostram parte da complexidade de modelar dados longitudinais \cite{braga2022random}, somando as dificuldades da ausência de normalidade à estruturação hierarquizada dos dados, motivo pelo qual o autores \cite{azzalini1999statistical, azzalini2003log, sahu2003new, sahu2004multivariate, braga2022random} têm se debruçado sobre esses campos na proposta de modelos cada vez abrangentes, contemplando distribuição assimétricas, bimodais, além da gaussiana (normal).  

% Aspectos como estes demonstram a complexidade na modelagem para uma fatia significativa dos dados longitudinais \cite{braga2022random}, motivo pelo qual estudos como os de  \cite{azzalini1999statistical, azzalini2003log, sahu2003new, sahu2004multivariate} têm despontado na proposição de novas tecnologias com a capacidade de ajustar-se cada vez melhor a estes tipos de dados, sobretudo nos campos das assimetrias e bimodalidades.

Considerando os desafios da distribuição dos dados, tendo base os estudos de \cite{azzalini1999statistical, azzalini2003log, sahu2003new, sahu2004multivariate}, o autor \cite{adolf2021} criou um novo modelo de regressão chamado Odd Log Logística Skew t-Student (OLLST) com a proposta de ajustar-se melhor na presença de assimetrias (à esquerda e à direita) e assimetrias de bimodalidades.   

Na aplicação do modelo OLLST, \cite{adolf2021} pôde constatar o bom desempenho dele na análise para diferentes tipos de distribuição, recomendando-o como alternativa para delineamentos em que a variável resposta não segue distribuição normal.

% Seguindo esta linha, Adolf Fernandes na sua dissertação de mestrado em 2021 \cite{adolf2021}, propôs um novo modelo o qual chamou Odd Log Logistic Skew t-Student (OLLST). Segundo ele, o modelo tem a capacidade de ajustar-se bem à assimetrias de ambos os lados, assimetria de bimodalidade à esquerda, à direita, além do comportamento normal. Tendo apresentado bons resultados na aplicação real para modelagem de dados da COVID-19, o modelo ainda não pode ser explorado na sua totalidade quanto a análise de dados longitudinais, mesmo que em teoria, tenha os pressupostos necessários para tal.

Tendo como base os resultados propostos por \cite{adolf2021} para análise de assimetrias e bimodalidades, esta dissertação analisou o comportamento do modelo OLLST no delineamento de dados reais, comparando-o ao modelo normal (NO), a fim de saber qual explicaria melhor a relação imposta entre as variáveis no ajuste de dados longitudinais. 

Durante o estudo sobre o modelo, verificou-se que, por pertencer à classe de modelos aditivos generalizados para posição, escala e forma (GAMLSS), a interpretação da saída de dados da biblioteca $\mathcal{R}$ para essa família de modelos tende a ser confusa, pouco intuitiva e com aumento da complexidade na proporção em que cresce o número de indivíduos. Por esse motivo, uma proposta de automação e adaptação da saída foi implementada com a possibilidade de uso ou não de uma interface gráfica de manipulação.  



\section{Organização}
\label{cap1:secOrganizacao}

Esta dissertação segue estruturada conforme o disposto. No Capítulo \ref{cap2:title} será apresentada uma revisão de conceitos relacionados a pesquisa com intuito de nortear o leitor quanto ao contexto do trabalho e tecnologias usadas para construção do modelo OLLST pelo autor \cite{adolf2021}. Nesse ponto, a proposta de medida repetida no tempo é incorporada ao ambiente da análise de regressão de modo que fiquem claro os desafios de inerentes à estas modelagens, suscitando uma abordagem mais eficiente.

No Capítulo \ref{cap3:title} serão discutidas as estrutura matemáticas de geração do modelo OLLST, sua incorporação à família de distribuição GAMLSS, a postulação do modelo hierárquico para modelagem de dados longitudinais e os mecanismos de estimação usado pela função. Será discorrido também sobre o momento de pesquisa que levou a proposta de criação para melhoria na análise e saída dos dados ajustados pelos modelos de regressão hierárquico.

No Capítulo \ref{cap4:title} são apresentados inicialmente os resultados obtidos para os ajustes de seis conjuntos de dados para o teor de carbono no solo de diferentes tipos de tratamentos em função do tempo, além da análise comparativa do modelo OLLST em razão do Normal. Neste capítulo também será apresentada a ferramenta desenvolvida com objetivo de facilitar a interpretação dos dados para ajustes longitudinais, com interface gráfica e rotinas de automação para comparações. 

Finalmente no Capítulo \ref{cap5:title} uma síntese da pesquisa será exposta com as considerações que culminaram neste estudo e um panorama dos resultados alcançados para duas frentes de trabalho proposta. 

 
\chapter{Conceitos}
\label{cap2:title}



\section{Regressão}
\label{cap2:secRegressao}

De acordo \cite{sahu2003new}, a análise de regressão é um método importante na estatística e tem por proposta conhecer os efeitos que algumas variáveis podem ou não exercer sobre outras. Ainda que não existe de fato uma relação casual entre as variáveis, com análise de regressão é possível relacioná-las por meio de uma expressão matemática, podendo desta forma estimar o valor de uma variável com base nos valores das demais. 

Tais relações podem ser genericamente representada por 

\begin{equation}
	\label{eqn21modregre}
	y=f(x_{1},x_{2},...,x_{k}),
\end{equation} em que $y$ é a variável dependente e os $x_{k}$ são as variáveis explicativas. Em um nível mais interno de conceitos os modelos regressão podem apresentar comportamento diferentes para estimação de valores, sendo eles paramétricas, não-paramétricas e semi-paramétricas.  

De maneira geral, o conceito e aplicação dos tipos de regressão está ligados diretamente ao conhecimento ou não da função $f$.  No primeiro caso, tem-se o conhecimento da função de distribuição $f$, no entanto, não se sabe o valor dos parâmetros. Na regressão não-paramétrica tem-se o comportamento inverso ao paramétrico, ao passo que a regressão semi-paramétrica tenta contemplar componentes paramétricos e não-paramétricos. 

\section{Medida repetida no tempo}
\label{cap2:secMedidaRepedida}

O estudo de medidas repetidas no tempo apresenta ampla aplicação, contemplado áreas como medicina, agronomia, economia, entre outras, tendo como principal característica a observação repetida de um indivíduo ao logo do tempo. Um exemplo disso ocorre quando de deseja modelar a variação diária da pressão sanguínea de indivíduos com diferentes tipos de quadros clínicos \cite{singer2000analysis}. 

Neste cenário, as aferições repetidas ao longo do tempo para o mesmo indivíduo estabelecem uma relação de associação entre o tempo e a variável resposta. Se os indivíduos deste cenário fossem submetidos a tratamentos medicamentosos diferentes e as aferições repetidas avaliasse a reação da droga, a análise comparativa entre indivíduos para cada tratamento deveria então considerar o tempo como fator de influência e assim concluir, por exemplo, qual o medicamento mais eficaz. 

% Existe hoje uma onipresença de estudos com dados longitudinais em que pesquisadores, com pouco ou nenhum contato com estatística, se deparam com diversos desafios ao analisá-los. Seja na saúde, agricultura, ciências sociais, biologia, educação e outros, o trabalho com dados correlacionados expande os horizontes de aplicações e extrapola os limites da estatística levando soluções e desafios a outra áreas.

% De maneira didática, a literatura apresenta uma gama de conceitos para o termo genérico "dados correlacionados", podendo também ser encontrado termos como: observações multivariadas, dados agrupados, medida repetida, dados longitudinais e dados espacialmente correlacionados. Em todos os casos, a característica essencial de estudos dessa natureza é repetição das observações ao logo do tempo \cite{verbeke1997linear}. Para este trabalho, considerou-se então o tempo \textbf{medida repetida no tempo}. 

No aspecto estrutural, os dados longitudinais seguem de modo que o vetor com $m_{i}$ respostas para o $i$-ésimo indivíduo $i=(1,...,n)$ qualifica o perfil individual da resposta, variando ainda no tempo $k=(1,...,m_{i})$, ou seja, cada indivíduo $i$ terá no tempo $k$ um vetor com $n$ variáveis resposta. Desse modo, é possível descrever o perfil de uma resposta $\mathbf{y}$ no tempo como $\mathbf{y}_{i}=(y_{i1}, y_{i2},...,y_{im_{i}})$ \cite{kac2007epidemiologia, singer2000analysis}

%Em medida repetida no tempo, segue que $N$ indivíduos são observados em $n_{j}$ ocasiões do tempo $j=(1,2,...n_{i})$ para à variável resposta $y$ e um conjunto de $p$ variáveis explicativas, de modo que $n_{i}$ é o número de observações do indivíduo $i$ e $n=\sum_{i=1}^{j_{N}}n_{i}$ é o número total de observações. Desse modo, $y_{i}^{\top }=(y_{i1},y_{i2},...,y_{im_{i}})$ é a representação do vetor de observações da variável resposta $y$ para indivíduo $i$. Ou seja, a estrutura básica composta que compõe um estudo de medida repetida no tempo é dada por um conjunto de variáveis, por um grupo de indivíduos e pelas medidas tomadas ao longo tempo para um mesmo indivíduo \cite{kac2007epidemiologia}. 

Este conceito é melhor compreendido ao ser observar a Tabela \ref{tab1cap2disposicao} em que é proposto um paradigma para organização de dados longitudinais, demonstrando os pontos de influência do tempo sobre as variáveis. Considerando $y_{ik}$ elemento do vetor de $m_{i}$ respostas no tempo $k$ para o $i$-ésimo individuo, além das variáveis explicativas $(X,W,V e Z)$, tem-se que o tempo está associado, por exemplo, as variáveis $y$ e $Z$, indicando uma relação de dependência entre eles.


\begin{table}[h]
    \begin{center}
        \caption{Estrutura para disposição de dados longitudinais \cite{singer2000analysis}.}
        \label{tab1cap2disposicao}\vspace*{0.3cm}
        \begin{tabular}{@{}lcccccc@{}}
            \toprule
             						& 				&				& \multicolumn{4}{c}{Covariáveis} \\ \cmidrule(l){4-7} 		
           	Indivíduos \tiny{$(i=1,...,n)$}		& Resposta \tiny{$m_{i}$} & Tempo \tiny{$(k=1,...,m_{i})$}			& X					& W				& V				& Z			\\ \midrule
           	1						& $y_{11}$		& $t_{11}$		& $x_{1}$			& $w_{1}$		& $v_{1}$		& $z_{11}$		\\		
           	1						& $y_{12}$		& $t_{12}$		& $x_{1}$			& $w_{1}$		& $v_{1}$		& $z_{12}$		\\	
           	.						& --			& --			& --				& --			& --			& --			\\
           	1						& $y_{1m_{1}}$	& $t_{1m_{1}}$	& $x_{1}$			& $w_{1}$		& $v_{1}$		& $z_{1m_{1}}$	\\	
           	2						& $y_{21}$		& $t_{21}$		& $x_{2}$			& $w_{2}$		& $v_{2}$		& $z_{21}$		\\	
           	2						& $y_{22}$		& $t_{22}$		& $x_{2}$			& $w_{2}$		& $v_{2}$		& $z_{22}$		\\	
           	.						& --			& --			& --				& --			& --			& --			\\
           	2						& $y_{2m_{2}}$	& $t_{2m_{2}}$	& $x_{2}$			& $w_{2}$		& $v_{2}$		& $z_{2m_{2}}$	\\	
           	n						& $y_{n1}$		& $t_{n1}$		& $x_{n}$			& $w_{n}$		& $v_{n}$		& $z_{n1}$		\\	
          	n						& $y_{n2}$		& $t_{n2}$		& $x_{n}$			& $w_{n}$		& $v_{n}$		& $z_{n2}$		\\	
            .						& --			& --			& --				& --			& --			& --			\\
            n						& $y_{nm_{n}}$	& $t_{nm_{n}}$	& $x_{n}$			& $w_{n}$		& $v_{n}$		& $z_{nm_{n}}$	\\	
        \end{tabular}
    \end{center}
\end{table}


Outro modo interpretar o comportamento primário de medida repetida no tempo se dá pela análise do gráfico de perfis individuais  (veja as Figuras \ref{fig1graficoPerfis1} e \ref{fig1graficoPerfis}), sendo possível inferir a variação na resposta $y$ ao longo do tempo, bem como em seu ponto inicial, justificando, por exemplo, o uso de efeito aleatório no intercepto durante o delineamento, a fim de oferecer uma margem maior de variabilidade ao modelo. 

 %Essa influência do tempo sobre as variáveis respostas pode ser melhor compreendida ao se plotar um gráfico de perfis das unidades amostrais em função do tempo, como mostra a figura \ref{fig1graficoPerfis}.

% \begin{figure}[h]
% 	\centering
% 	\includegraphics[width=0.8\linewidth]{perfil_individual_medio_tratamento_separado}
% 	\caption{Gráfico de perfis com perfil individual médio por individuo}
% 	\label{fig1graficoPerfis1}
% \end{figure}

% \begin{figure}[h]
%   \centering
%   \includegraphics[width=0.5\linewidth]{perfil_individual_medio_tratamento}
%   \caption{Gráfico de perfis com perfil individual médio em conjunto}
%   \label{fig1graficoPerfis}
% \end{figure}

Entre os modelos dispostos na literatura os mais difundidos no delineamento de dados longitudinais trabalham com a pressuposição de normalidade da variável resposta ou dos erros \cite{singer2000analysis}. Este é o caso, por exemplo, dos modelos lineares misto (ou modelo de efeito aleatório), que apesar da amplitude na modelagem ao incorporar o efeito aleatório, apresentam limitações no campo da distribuição. 

\section{Modelo linear de efeito misto}
\label{cap2:secModeloLinearMisto}

Proposto por \cite{laird1982random}, o modelo de regressão linear é também utilizando como uma das técnicas de análise de dados longitudinais. Dado por 
\begin{equation}
	\label{eqn22modelinearmisto}
	y_{i}=X_{i}\beta +Z_{i}b_{i}+e_{i}
\end{equation}, em que $y_{i}$ é o vetor de resposta de dimensão $n_{i} i=1,2,...,m$ e $Z_{i}$ são matrizes de delineamento (ou incidência) conhecidas, de dimensões $n_{i}\times p$e $n_{i}\times q$ respectivamente, $\beta$ é um vetor $p$-dimensional contendo os efeitos fixo, $b_{i}$ é um vetor $q$-dimensional contendo os efeitos aleatórios, e $e_{i}$ é o vetor de resíduos de dimensão $n_{i}$. Diferente dos modelos padrões de regressão linear, este incorpora ao modelo a dependência das observações e a estrutura de correlação dos erros \cite{fausto2008modelo, ferreira2012analise, freitas2005alternativas, pereira2013eficiencia}.

Por outro lado, a distribuição dos dados pode impactar diretamente na eficiência do modelo, pois exige-se uma pressuposição de normalidade \cite{azzalini1999statistical}, uma premissa falha caso o conjunto de dados possa apresentar um comportamento foram do padrão esperado. Na mesma linha, ao se trabalhar com medida repetida no tempo, tem-se o fato de que a distribuição da resposta, na maioria dos casos, apresenta uma zona de assimetria em direção aos maiores tempos, o que tornaria inapropriado o uso do modelo \ref{eqn22modelinearmisto} \cite{colosimo2006analise}. 

A figura \ref{fig22histno}, por exemplo, apresenta um histograma para determinado conjunto de dados cuja a distribuição não é normal, ao passo que a \ref{fig23histno} mostra um conjunto com distribuição esperado pelos modelos lineares. Em ambos os casos, a linha vermelha sob o gráfico mostra a função de densidade e probabilidade (fdp) da normal, mostrando as limitações do modelo quando se percebe de dissonâncias na distribuição. 

% \begin{figure}[h]
% 	\centering
% 	\includegraphics[width=0.6\linewidth]{hist_no_a1}
% 	\caption{Histograma para dados sem distribuição normal.}
% 	\label{fig22histno}
% \end{figure} 

% \begin{figure}[h]
% 	\centering
% 	\includegraphics[width=0.6\linewidth]{hist_no_a2}
% 	\caption{Histograma para dados com distribuição normal.}
% 	\label{fig23histno}
% \end{figure} 

Grande parte dos estudos com medida repetida no tempo enfrentam as limitações da modelagem do modelo linear transformando a variável resposta para tentar retornar ao modelo linear-normal, quando não, acabam utilizando um componente sistemático não-linear nos parâmetros e uma distribuição assimétrica para o componente estocástico. Em ambos os casos, o pesquisador pode se deparar com outras distribuições assimétricas ou mesmo bimodais para o erro, que não sejam possíveis de retornar para o modelo linear \cite{colosimo2006analise}, exigindo modelos mais flexíveis para o ajuste. 

Deve-se considerar, neste contexto, o momento histórico para o uso de cada modelo de regressão. Criado em meados de 1964, devido as limitações computacionais para cálculos mais complexos ou de grande volumes de dados, a transformação Box-Cox \cite{box1964analysis}, de maneira minimalista, produzia uma normalidade aproximada para o ajuste dos modelos da época. Neste mesmo período, métodos como da máxima verossimilhança que envolviam cálculos mais complicados eram simplesmente impraticáveis. Contudo, hoje, com a melhora na performance dos computadores, o estudo de métodos simplista está sendo desistimulado em detrimento a técnicas mais eficientes \cite{maddala2003introduccao}.

\section{Modelos aditivos}
\label{cap2:secModelosAditivos}

Os modelos aditivos são uma generalização do modelo linear tradicional, por tanto, acompanha a importante característica de verificar a contribuição que cada variável exerce sobre a variável de interesse \cite{de2001modelos}. Tomando como base a equação \ref{eqn221regparametrica}, considerando $n$ pares de observações $(x_{i},y_{i}), i=(1,2,...,n)$ e sendo $f$ a função que estabelece a relação entre as variáveis $X$ e $Y$ da forma

\begin{equation}
	\label{eqn23modeloaditivo1}
	y_{i}=f(x_{i})+e_{i}.
\end{equation} 

Considerando um conjunto de $k$ variáveis explicativas representadas em uma matriz $X$, de dimensão $n\times k$ com a $i$-ésima linha dada por $X_{i} = (x_{i1}, x_{i2},...,x_{ik})$, teremos $f$ de modo que $y_{i}=f(x_{i1}, x_{i2},...,x_{ik})+e_{i}$.

Os autores \cite{buja1989linear, de2010engenharia}, desenharam uma nova família de modelos, chamados modelos aditivos, de modo que fosse possível também aplicar a regressão linear no efeito as variáveis regressoras, de modo que modelo proposto na equação \ref{eqn23modeloaditivo1} passou a ser 

\begin{equation}
	\label{eqn23modeloaditivo2}
	y_{i}=f_{1}(x_{i1})f_{2}(x_{i2})+...+f_{k}(x_{ik})+e_{i},
\end{equation}

A principal vantagem desta proposta seria a superação do problema conhecido na literatura como "maldição da dimensionalidade" \cite{koppen2000curse}, já que cada função $f$ do somatório é estimada de modo univariado. 

Apesar da melhora nas abordagens dos modelos existentes para análise de dados longitudinais (ver por exemplo \cite{buja1989linear}, \cite{de2010engenharia} \cite{sahu2003new} e \cite{azzalini2003log}), ainda existem lacunas no estudo de modelos assimétricos e bimodais, visto que fenômenos desta natureza acrescentam uma camada extra de complexidade no fator de identificação do comportamento da população. Na bimodalidade, por exemplo, é comum se deparar com a miscigenação de duas populações, em que a zona de densidade e probabilidade dos dados caracteriza-se por duas modas, aumentado consideravelmente a complexidade de se inferir comportamento dos indivíduos \cite{wang2009bimodality}.



\section{Modelos Aditivos Generalizados de Localização, Escala e Forma (GAMLSS)}
\label{cap2:secGAMLSS}

A construção da família de modelos GAMLSS é resultado de uma linha do tempo evolutiva dos modelos tradicionais e suas limitações somado ao avanço da capacidade computacional para resolução de problemas complexos. 

Como parte deste cenário, os modelos lineares generalizados (\textit{Generalized Linear Models} - GLM) \cite{nelder1972generalized} assim como os modelos aditivos generalizados (\textit{Generalized Additive Models} - GAM) \cite{buja1989linear}, derivado dos lineares tradicionais, por muito tempo ocuparam lugar de destaque na literatura como referência na análise de regressão univariada. 

%Estes modelos assumem que a distribuição da variável resposta pertencem a família exponencial, modelando a sua média $\mu$ a partir das variáveis explicativas, limitando-se ainda a comportamentos gaussianos da distribuição.

Tanto nos modelos GLM quanto nos GAM, a pressuposição de que a variável resposta \textbf{$y$} tenha distribuição pertencente a família exponencial, implica na modelagem da média \textbf{$\mu$} em função das variáveis preditoras. Nesses modelos a variância de \textbf{\text{$y$}} depende de um parâmetro de dispersão constante \textbf{$\phi$} e da média \textbf{$\mu$},  ou seja, a variância, assim como a assimetria e a curtose não são modeladas explicitamente em termos de variáveis preditoras e sim por meio da dependência da média \textbf{$\mu$} \cite{rigby1996madam}.

Estudos como \cite{breslow1993approximate, breslow1995bias, nelder1972generalized}, apresentam uma nova classe de modelo como resultado de um \textit{crossover} (cruzamento de estilos) entre o modelo linear generalizado e modelos linear misto com o acréscimo de um termo (quase sempre normal). Os derivados desta nova família foram chamados de modelos generalizado misto (\textit{Generalizerd Linear Mixed Model} - GLMM) e mesmo apresentando maior flexibilidade que os GLM e GAM, a pressuposição de que $y$ tem sua distribuição pertencente a família exponencial torna-o limitante em alguns aspectos.

Como resultado destas barreiras, os autores \cite{rigby2005generalized} propuseram uma nova abordagem para análise de relações entre as variáveis. Nesse modelo de regressão (semi)paramétrica, seu comportamento paramétrico faz referência a necessidade de uma suposição quanto distribuição para variável resposta, ao passo que o comportamento "semi" diz respeito a amplitude na modelagem dos parâmetros de distribuições e funções para formas lineares, não lineares e suavizações de variáveis explicativas.

No GAMLSS, diferente do GLM e GAM citados anteriormente, a variável resposta $y$ não precisa pertencer à família exponencial. Neste caso, a distribuição de $y$ passa a pertencer à uma família mais genérica ($\mathcal{D}$) dada por $D(y|\mu,\sigma,\nu,\tau)$, em que $D\in \mathcal{D}$, podendo ser qualquer distribuição, seja ela contínua com assimetrias ou curtose acentuadas ou mesmo discreta. Uma outra vantagem dos modelos derivados do GAMLSS é a possibilidade de modelagem de todos os parâmetros da distribuição condicional de $y$ e não apenas a média $\mu$. \cite{rigby1996madam}. Somando a isso, existe ainda o aspecto que refere-se a facilidade de acesso e manuseio do modelo computacionalmente desenhado e implementado. A biblioteca \texttt{gamlss} da linguagem de programação $\mathcal{R}$ permite ajustar mais de 50 distribuições diferentes, já incorporadas no pacote, podendo ainda ajustar novos modelos que atendam as premissas do GAMLSS.

\subsection{Definição do modelo}
\label{cap221:definicaomodelogamlss}

Em sua estrutura, os $p$ parâmetros $\theta ^{\top}=(\theta_{1},\theta_{2},...,\theta_{p})$ de uma função de densidade e probabilidade $f(y|\theta)$ são modelados utilizando termos aditivos, presumindo que para $i=(1,2,...,n)$ as observações $y_{i}$ são independentes e condicionais a $\theta^{i}$. Neste caso, a função de densidade e probabilidade é dada por $f(y_{i}|\theta^{i})$, em que $\theta^{i\top}=(\theta_{i1},\theta_{i1},...,\theta_{ip})$ é um vetor de $p$ parâmetros relacionado às variáveis explanatórias e efeitos aleatórios. De modo que o modelo GAMLSS pode ser descrito por

\begin{equation}
	\label{eqn22cap2gamlss1}
	g_{k}(\theta _{k})=\eta_{k}=X_{k}\beta _{k}+\sum_{j=1}^{j_{k}}Z_{jk}\gamma _{jk}, 
\end{equation}, em que $\theta_{k}$ e $\eta_{k}$ são vetores de parâmetros $n\times 1$, $X_{k}$ e $Z_{jk}$ são covariáveis fixas, conhecidas e de ordem $n\times J_{k}^{'}$ e $n\times q_{jk}$, respectivamente, ao passo que $\gamma _{jk}$ é uma variável aleatória $q_{jk}$-dimensional.

Grande parte dos estudos práticos para análise de regressão utilizam um máximo de quatro parâmetros $(p=4)$, usualmente caracterizados pela posição $(\mu)$, escala $(\sigma)$, assimetria $(\nu)$ e curtose $(\tau)$, sendo $\theta _{1} = \mu$ e $\theta _{2} = \sigma $ parâmetros de posição (ou locação) e $\theta _{3} = \mu$ e $\theta _{3} = \sigma $ parâmetros de forma. Com isto, tem-se os seguintes modelos:

\begin{equation}
\label{eqn11cap1}
	g_{k}(\theta _{k})=\eta_{k}=X_{k}\beta _{k}+\sum_{j=1}^{j_{k}}Z_{jk}\gamma _{jk}
\end{equation}

\begin{equation}
\label{eqn12cap1}
	g_{1}(\mu)=\eta_{1}=X_{1}\beta _{1}+\sum_{j=1}^{j_{1}}Z_{j1}\gamma _{j1}
\end{equation}

\begin{equation}
\label{eqn13cap1}
	g_{2}(\sigma)=\eta_{2}=X_{2}\beta _{2}+\sum_{j=1}^{j_{2}}Z_{j2}\gamma _{j2}
\end{equation}

\begin{equation}
\label{eqn14cap1}
	g_{3}(\nu)=\eta_{3}=X_{3}\beta _{3}+\sum_{j=1}^{j_{3}}Z_{j3}\gamma _{j3}
\end{equation}

\begin{equation}
\label{eqn15cap12}
	g_{4}(\tau)=\eta_{4}=X_{4}\beta _{4}+\sum_{j=1}^{j_{4}}Z_{j4}\gamma _{j4}
\end{equation}

\subsection{Família GAMLSS}
\label{cap251:familiagamlss}

Considerando que função de densidade e probabilidade $f(y|\theta)$ do modelo \ref{eqn22cap2gamlss1} pode pertencer à $\mathcal{D}$, dado por $D(y|\mu,\sigma,\nu,\tau)$, em que $D\in \mathcal{D}$, ao se ajustar um modelo GAMLSS utilizando a biblioteca \texttt{gamlss} da linguagem de programação $\mathcal{R}$ é necessário atentar-se unicamente para que a função $f(y|\theta)$ assim como sua primeira derivada, para cada parâmetro $\theta$, sejam calculáveis.

As tabelas \ref{tab2cap2discrete} e \ref{tab3cap2continuous} apresenta uma lista de distribuições já incorporadas à biblioteca \texttt{gamlss} do $\mathcal{R}$.

\begin{table}[h]
	\begin{center}
		\caption{Exemplos de distribuições discretas \cite{rigby2005generalized}.}
		\label{tab2cap2discrete}\vspace*{0.3cm}
		\begin{tabular}{@{}llcccc@{}}
			\toprule
			&          & \multicolumn{4}{c}{Função de Ligação}  \\ \cmidrule(l){2-6} 
			Distribuição                              	& Nome   & $\mu$ & $\sigma$ & $\nu$ & $\tau$ \\ \midrule
			Beta Binomial								& BB()			& Logit		& Log			& --				& --\\
			Binomial									& BI()			& Logit		& --			& --				& --\\
			Delaporte									& DEL()			& Log		& Log			& Logit				& --\\
			Negative Binomial type I					& NBI()			& Log		& Log			& --				& --\\
			Negative Binomial type II					& NBII()		& Log		& Log			& --				& --\\
			Poisson										& PO()			& Log		& --			& --				& --\\
			Poisson Inverse Gaussian					& PIG()			& Log		& Log			& --				& --\\
			Sichel										& SI()			& Log		& Log			& Identity			& --\\
			Sichel ($\mu$ the mean)						& SICHEL()		& Log		& Log			& Identity			& --\\
			Zero Inflated Poisson						& ZIP()			& Log		& Logit			& --				& --\\		\bottomrule
		\end{tabular}
	\end{center}
\end{table}
\pagebreak
\begin{table}[]
	\begin{center}
		\caption{Exemplos de distribuições contínuas \cite{rigby2005generalized}.}
		\label{tab3cap2continuous}\vspace*{0.3cm}
		\begin{tabular}{@{}llcccc@{}}
			\toprule
			&          & \multicolumn{4}{c}{Função de Ligação}  \\ \cmidrule(l){2-6} 
			Distribuição                              & None   & $\mu$ & $\sigma$ & $\nu$ & $\tau$ \\ \midrule
			Beta                                       & BE()     & Logit              & Logit                 & --                 & --                  \\
			Beta Inflated (at 0)                       & BEOI()   & Logit              & Log                   & Logit              & --                  \\
			Beta Inflated (at 1)                       & BEZI()   & Logit              & Log                   & Logit              & --                  \\
			Beta Inflated (at 0 and 1)                 & BEINF()  & Logit              & Logit                 & Log                & Log                 \\
			Box-Cox (Cole \& Green)                    & BCCG()   & Identity           & Log                   & Identity           & --                  \\
			Box-Cox Power Exponential                  & BCPE()   & Identity           & Log                   & Identity           & Log                 \\
			Box-Cox \textit{t}                         & BCT()    & Identity           & Log                   & Identity           & Log                 \\
			Exponential                                & EXP()    & Log                & --                    & --                 & --                  \\
			Exponential Gaussian                       & exGAUS() & Identity           & Log                   & Log                & --                  \\
			Exponential Gen. Beta type II              & EGB2()   & Identity           & Identity              & Log                & Log                 \\\bottomrule
		\end{tabular}
	\end{center}
\end{table}

A notação para descrever o uso dos modelos acima em $\mathcal{D}$ é dada por:

\begin{equation}
	\label{eqn15cap1}
	y \sim \mathcal{D} \left\{g_1(\theta_1) = t_1, g_2(\theta_2) = t_2, ..., g_p(\theta_p) = t_p\right\},
\end{equation}, em que $\theta_{1},...,\theta_{p}$ são os parâmetros de $\mathcal{D}$, $g_{1},...,g_{p}$ são as funções de ligação e $t_{1},...,t_{p}$ são as fórmulas dos modelos para os temos explanatórios e/ou efeitos aleatórios nos preditores $\eta_{1},...,\eta_{p}$, respectivamente.

\section{Observações finais}
\label{cap2:observacoesFinais}

Neste capítulo, foram apresentados os conceitos e tecnologias que subsidiaram a construção do modelo OLLST de modo que, para o leitor, seja possível compreender os conceitos tratados no decorrer da pesquisa. Este trabalho tem como foco o estudo do modelo OLLST na análise de dados longitudinais, para tal, é necessário entender os conceitos que permeiam este cenário, tais como: analise/modelo de regressão, medida repetida no tempo, entre outros.

Sob esse panorama, os modelo Modelos Aditivos Generalizados de Localização, Escala e Forma (GAMLSS) tem fundamental importância, pois serviam de base para construção da distribuição descrita no capítulo \ref{cap3:title}, de modo que os conceitos apresentados até agora se culminem em uma nova pesquisa.   
\chapter{Materiais e Métodos}
\label{cap3:title}

\section{Modelo de regressão OLLST}

O modelo Odd Log-logística Skew $t$-Student, indicado neste trabalho por OLLST, é resultado da extensão da ST, de modo que, dada a função $W[G(x)]$, onde $g(x)$ e $G(x)$ é a função densidade e distribuição de densidade, respectivamente, da Skew $t$-Student descrita por:

\begin{equation}
	\label{eqn1cap3}
	W[F(x)]=\frac{G(x, \mu; \sigma)}{1-\bar{G}(x,\mu; \sigma)},
\end{equation} foi possível, baseado na transformação de $T-X$, propor um novo modelo para família de distribuições GAMLSS cuja a função de distribuição é dada por: 

\begin{equation}
	\label{eqn28cap3}
	F(x;\mu, \sigma, \lambda, \alpha, \nu) =\int_{0}^{W[G(x)]} \frac{\alpha \times s^{\alpha - 1}}{(\alpha - s^\alpha)^2} ds = \frac{G(x)^\alpha} {G(x)^\alpha + [1 - G(x)]^\alpha}.
\end{equation}

A função de densidade de probabilidade é descrita como:

\begin{equation}
	\label{eqn29cap31}
	f(x;\mu, \sigma, \lambda, \alpha, \nu) = \frac{\alpha \times G(x)^{\alpha-1} g(x)[1 - G(x)]^{\alpha-1}}{\{G(x)^\alpha + [1 - G(x)]^\alpha\}^2}.
\end{equation}

A funções própria para o parâmetro \textbf{$\alpha$ }é definida por:

\begin{equation*}
	\label{eqn299cap31}
	\alpha = \frac{log\left(\frac{F(x)}{1 - F(x)}\right)}{log\left(\frac{G(x)}{1 - G(x)}\right)}.
\end{equation*} tendo ainda a função quantílica, dada por $Q(\alpha, \mu)$ em função de $x$ representada como segue:
\begin{equation}
	\label{eqn30cap31}
	Q(\alpha, u) = G(x)^{-1}\left[\frac{u^{\frac{1}{\alpha}}}{(1 - u)^{\frac{1}{\alpha}}+u^{\frac{1}{\alpha}}}\right].
\end{equation}


A flexibilidade deste novo modelo reside na restrição dos parâmetros $\alpha$ e $\lambda$, e consequente variação de $\nu$, de modo que:
\begin{enumerate}
	\item Quando $\alpha \neq 1$ e $\lambda \neq 0$:
	\begin{itemize}
		\item $\nu = 1$ (OLLST Cauchy);
		\item $\nu = 4 - 20$ (OLLST);
		\item $\nu \geq 20$ (OLLST).
	\end{itemize}
	
	\item Quando $\alpha = 1$ e $\lambda \neq 0$:
	\begin{itemize}
		\item $\nu = 1$ (ST Cauchy);
		\item $\nu = 4 - 20$ (ST);
		\item $\nu \geq 20$ (SN).
	\end{itemize}
	
	\item Quando $\alpha = 1$ e $\lambda = 0$:
	\begin{itemize}
		\item $\nu = 1$ ($t$-Student Cauchy);
		\item $\nu = 4 - 20$ ($t$-Student);
		\item $\nu \geq 20$ (Normal).
	\end{itemize}
\end{enumerate}

De maneira visual, a figura \ref{fig7differentGraphs} apresenta a variação do comportamento da função de densidade OLLST na medida que os parâmetros mudam seu valor. Em ambas as imagens, os valores para $\mu$ e $\sigma$ foram fixados em $0$ e $1$, respectivamente. Enquanto na figura \ref{fig7differentGraphs}(a) $\lambda = 2$ e $\nu=13$, com variação no valor de $\alpha$, conforme a legenda, na figura \ref{fig7differentGraphs}(b) $\alpha$ foi fixado em $0,2$ e $\nu=13$ variando somente o valor para $\lambda$.

% \begin{figure}[!htb]
% 	\begin{center}
% 		\hspace{0.2cm}(a)\hspace{5cm} (b)\\
% 		\includegraphics[width=7cm,height=4.7cm]{fig12}
% 		~~\includegraphics[width=7cm,height=4.7cm]{fig13}\\
% 		\caption{Variação da densidade do OLLST.}
% 		\label{fig7differentGraphs}
% 	\end{center}
% \end{figure}

O modelo foi codificado em $\mathcal{R}$ podendo ser utilizado pela biblioteca \texttt{gamlss} passando o valor "OLLST" como parâmetro de escolha da distribuição \textit{(Veja o apêndice \ref{appendixA})}.

A estimação dos parâmetros do modelo fica a cargo do método de estimação de máxima verossimilhança, mostrados na seção \ref{cap5:secEstimacao}, visto que não há necessidade de uma suposição, ou quando há é mínima, da função de densidade e probabilidade. Além disso, o modelo considera ainda o uso de quatro parâmetros, com a justificativa de ser amplo o suficiente para descrever os fenômenos no mundo real e possibilidade de ajustar-se mais adequadamente à assimetria e bimodalidade à direita e/ou à esquerda \cite{adolf2021}. 

\subsection{Estimação}
\label{cap5:secEstimacao}

A estimação de parâmetros para distribuição OLLST com efeitos aleatórios para dados longitudinais é feita através do método da máxima verossimilhança, onde se maximiza o logaritmo da função de verossimilhança marginal, por meio da integral da função de verossimilhança em relação aos efeitos aleatórios $w_{i}$. No $i$-ésimo individuo o vetor da variável resposta pode ser expresso por $Yn_{i} = [Y_{i11},\ldots, Y_{i1K};Y_{i21}, \ldots,Y_{i2K};\ldots;Y_{in_{i}1},\ldots,Y_{in_{i}K}]^{T}$ e função de verossimilhança condicional dos efeitos aleatórios (independência dentro dos indivíduos) para esse individuo pode ter a seguinte forma:

\begin{equation}
	\label{eqn121cap3}
	L_{i}(y_{ijk}|w_i)=\prod_{j=1}^{n_i}\prod_{k=1}^{K}\,f(y_{ijk}|w_i),
\end{equation} onde o $f(y_{ijk}|w_i)$ é a função de densidade e dada por por (\ref{eqn399cap3}).

\begin{equation}
	\label{eqn399cap3}
	f(x;\mu, \sigma, \lambda, \alpha, \nu) = \frac{\alpha \times G(x)^{\alpha-1} g(x)[1 - G(x)]^{\alpha-1}}{\{G(x)^\alpha + [1 - G(x)]^\alpha\}^2}.
\end{equation}

Nesse contexto, considerando a independência das variáveis aleatórias $w_i$ e $Y_{ijk}$ a contribuição do $i$-ésimo individuo para a função de verossimilhança marginal é:

\begin{equation}
	\label{eqn122cap3}
	\int L_{i}(y_{ijk}|w_i)\,f(w_{i};\sigma_{\beta_{0}}^{2})d w_i,
\end{equation} sendo $f(w_{i};\sigma_{\beta_{0}}^{2})$ a densidade normal para efeitos aleatórios definidos pela equação \ref{eqn122cap3} e $L_{i}(y_{ijk}|w_i)$ é expressa por:

\begin{equation}
	\label{eqn123cap3}
	L(\theta ) = \frac{1}{\sigma _{\beta _{0}}\sqrt{2\pi }}    \prod_{j=1}^{n_{1}}   \prod_{k=1}^{K}  f(y_{ijk} | w_{i})   exp \left \{ \ -{\frac{1}{2}}{}   \left (   \frac{w_{i}}{\sigma _{\beta _{0}}}  \right )  \right  \} dw_{i},
\end{equation}

Considerando um conjunto de dados $(y_{11k}, \textrm{x}_{11k}),\ldots$ , $(y_{1n_ik},\textrm{x}_{1n_ik}),$, $\ldots$ , $(y_{m1k},\textrm{x}_{m1k})$, $\ldots$ , $(y_{mn_ik},\textrm{x}_{mn_ik})$ de $n$ observações onde $n=(n_1+\ldots+n_i, \textrm{x}_{ijk})$ é o vetor de covariáveis associados ao $i$-ésimo individuo, na $j$-ésima repetição de indivíduo e $k$-ésima  repetição no tempo, o logaritmo da função de verossimilhança marginal em (\ref{eqn123cap3}) pode ser descrito como (\ref{eqn124cap3}): 

\small{
	\begin{eqnarray}
		\label{eqn124cap3}
		l(\theta)&=&\sum_{i=1}^{m}\log\left\{\frac{ \hat\sigma2\sqrt{2}}{{\sigma_{{\beta_{0}}}  \sqrt  {\pi}
		}}\sum_{p=1}^{q}\,v_p^{+}\,\,\prod_{j=1}^{n_i}\,\prod_{k=1}^{K}\,
		\frac{\alpha\,\phi\left(z_{ijk}\right)\Phi\left( {\lambda z_{ijk}}
			\right)\Phi_{SN}^{\alpha-1}\left(z_{ij};\lambda\right)
			[1-\Phi_{SN}\left(z_{ijk};\lambda\right)]^{\alpha-1}}{\sigma\Big\{\Phi_{SN}^{\alpha}\left(z_{ijk};\lambda\right)
			+\left[1-\Phi_{SN}\left(z_{ijk};\lambda\right)\right]^{\alpha}\Big\}^{2}}\right.\nonumber\\
		&&
		\,\times\left.\exp  \left\{{-
			\frac{1}{2}\left({\frac{{s_p^{+}}}{\sigma_{\beta_{0}}}} \right)^2
		}\right\}\right\}
	\end{eqnarray}
}em que  $z_{ijk}=(y_{ijk}-\mu_{ijk})/{\sigma}$.

Vale ressaltar que existe uma dificuldade na estimação de parâmetros pelo método de máxima verossimilhança e esta reside na avaliação das integrais da função de verossimilhança. Neste método, ao se maximizar o $l(\theta)$ o pesquisador pode se deparar com $m$ integrais de efeitos aleatórios $w_{i}$ para variável, cuja a solução analítica é inviável na maioria dos casos. Em contrapartida, o uso de métodos de integração numérica pode representar a solução esperar para contornar está dificuldade. Para este trabalho, considerou-se o método de integração numérica por quadratura de gauss-hermite (Código $\mathcal{R}$ no apêndice \ref{cap2:secGaussHermite}).

\subsection{Quadratura de Gauss}
\label{cap1:secGauss}

A maioria dos métodos usados para calcular a integral de uma função utilizam técnicas de aproximação do valor da integral por meio um polinômio em uma ou mais regiões de interesse. A depender da complexidade da função, o cálculo pode ser fatorado em uma função de pesos a fim de entregar o valor que mais se aproxima do real \cite{golub1969calculation}.

As quadraturas gaussianas se enquadram nesse contexto, pois, apesar de apresentar intervalos irregulares e possíveis zonas de agrupamentos, também fazem uso de métodos por aproximação polinomial de grau crescente, onde os nós (nodos) $x_{i}, i=(1,2,3...,n)$ da quadratura, são as raízes do polinômios $P_{n}(x)$ para todo ponto menor que $n$ \cite{schwartz1996theory}.

A ideia principal por trás dos métodos da quadratura de gauss é substituir a soma integral, por uma soma discreta, da seguinte forma:
\\
\begin{equation}
	\label{eqn12cap2}
	\int_{b}^{a} f(x)dx \approx \sum_{i=1}^{n} \omega _{i}f(x_{i}),
\end{equation} onde \textit{n} é o número de pontos, \textit{F(x)} refere-se a função no ponto de integração, $f(x_{i})$ é o valor da coordenada para o ponto de integração e $\omega_{i}$ é a função peso, também para o ponto de integração \cite{da2017integraccao}.

Ao substituir a soma integral pela soma discreta na Quadratura de Gauss, cria-se uma função peso que tem por objetivo distribuir o valor aproximado e o valor exato da integral dentro do intervalo em análise, considerando, claro, que a função $F(x)$ seja conhecida. Na Quadratura de Laguerre (\ref{eqn13cap2}) de $n$ pontos, por exemplo, os nós da quadratura $x_{i}$ são as raízes do $n$-ésimo polinômio de Laguerre. \cite{barbosaestudo}

\begin{equation}
	\label{eqn13cap2}
	\int_{-1}^{1} f(x)dx = \omega_{1}f(x_{1}) + \omega_{2}f(x_{2}) + \omega_{3}f(x_{3})+...+\omega_{i}f(x_{i}) = \sum_{i=1}^{n}\omega_{i}f(x_{i})
\end{equation}

Entre os métodos de integração que se vale dos conceitos de quadratura de gauss, destacam-se:

\begin{itemize}
	\item Fourier;
	\item Legendre;
	\item Chebyshev;
	\item Laguerre; e
	\item Hermite \textit{(escolhida para uso e aplicação neste trabalho)}.
\end{itemize}


\subsection{Quadratura de Gauss-Hermite}
\label{cap2:secGaussHermite}

O método de integração numérica por Gauss-Hermite apresenta duas versões com bases similares, mas os conceitos e aplicações diferentes, sendo elas: não adaptativa e adaptativa. Seu conceito é a extensão da quadratura Gaussiana sendo expresso por \ref{eqnHermit11cap1}

\begin{equation}
	\label{eqnHermit11cap1}
	\int_{-\infty }^{+\infty }g(x)dx=\int_{-\infty }^{+\infty }exp\{-x^{2}\}dx\approx\sum_{k=1}^{q}v_{k}f(s_{k})
\end{equation}

onde $s_{k}$ são as raizes do polinômio de Hermite de grau $q$ e $v_{k}$ são os pesos associados aos números de pontos da quadratura e do polinômio de Hermite $H_{q-1}(x)$ avaliado em $s_{k}$.	

Nas condições em que $g(x)$ é um polinômio de grau $2q-1$ é possível obter o valor exato da integral por meio de $\sum_{k=1}^{q}v_{k}f(s_{k})$ o que leva ao entendimento que a aproximação do valor da integral está associado ao número de pontos da quadratura. 

A base da quadratura de Gauss-Hermite gira em torno da função de peso $v(x)$ que, na maioria dos cálculos de integração, pode ser representada pela função de densidade e probabilidade de uma variável aleatória, onde $g(x)$ é o produto de $v(x)$ com uma outra função $gf(x)$ \ref{eqnHermit12cap1} podendo $v(x)$ \ref{eqnHermit121cap1} se expressa como:  

\begin{equation}
	\label{eqnHermit12cap1}
	g(x) = v(x)f(x) ,
\end{equation}

\begin{equation}
	\label{eqnHermit121cap1}
	v(x) = \frac{1}{\sqrt{2\pi \alpha^{2} }}\exp \left \{ -\frac{(x-\mu )^{2}}{2\sigma ^{2}} \right \},
\end{equation} contexto em que, sendo $z = (x-\mu ) / \sqrt{2\sigma ^{2}}$, a identidade da função pode ser representada por: 

\begin{equation}
	\label{eqnHermit13cap1}
	\int_{-\infty }^{+\infty } g(x)dx=\int_{-\infty }^{+\infty }f(x)v(x)dx=\int_{-\infty }^{+\infty }\pi ^{-1/2}f(\sqrt{2}\sigma z+\mu )\exp {\left \{ -z^{2} \right \}}dz
\end{equation} e calculada por Gauss-Hermite da seguinte forma \ref{eqnHermit14cap1}: 

\begin{equation}
	\label{eqnHermit14cap1}
	\int_{-\infty }^{+\infty } g(x)dx=\int_{-\infty }^{+\infty }f(x)v(x)dx\approx \sum_{k=1}^{q}\frac{v_{k}}{\sqrt{\pi }}f(\sqrt{2}\sigma s_{k}+\mu )
\end{equation}

\textbf{Observação 2.2.1:} \textit{Em $\mathcal{R}$, existem funções derivados do pacote \textbf{GHQp} que montam os vetores de pesos e pontos para métodos de integração por gauss, aceitando como parâmetros o número de pontos e o método usado. (Código $\mathcal{R}$ no apêndice \ref{appendix:h}).}

Considerando \ref{eqnHermit18cap1}

\begin{equation}
	\label{eqnHermit18cap1}
	H_{q}(x) = (-1)^{q}\exp \left \{ x^{2} \right \}\frac{d^{q}}{dx^{q}}(\exp )\left \{ -x^{2} \right \}
\end{equation} como sendo o polinômio de Gauss-Hermite de grau $q$ e 

\begin{equation}
	\label{eqnHermit19cap1}
	v_{k}=\frac{2^{q+1}q!\sqrt{\pi }}{\left [ H'_{q} (x) \right ]^{2}}s
\end{equation} o vetor de pesos desse polinômio, ao aplicar em um exemplo da quadratura de Gauss-Hermite com polinômio de dois pontos, ou seja, $q=2$ na integral 

\begin{equation}
	\label{eqnHermit120cap1}
	\int_{-\infty }^{+\infty }\exp \left \{ -x^{2} \right \}x^{2}dx
\end{equation}, observa-se que: 


\begin{equation}
	\label{eqnHermit122cap1}
	\int_{-\infty }^{+\infty }\exp \left \{ -x^{2} \right \}x^{2}dx \approx \frac{\sqrt{\pi }}{2}\left [ f\left ( \frac{\sqrt{2}}{2} \right ) + f\left ( -\frac{\sqrt{2}}{2} \right )\right ],
\end{equation} desse modo $f(x)=x^{2}$, $H_{q}(x)=4x^{2}-1$, os zeros do polinômios são $x_{1}\models +\sqrt{2}/2 $  e $x_{2}= -\sqrt{2}/2 $ e $v_{1}=v_{1}=\frac{\sqrt{\pi }}{4(\sqrt{2}/2)^{2}}$, valor exato da integral \ref{eqnHermit120cap1} é $\frac{\sqrt{\pi }}{2}$, como mostra a tabela \ref{tableHermite11cap3} para polinômios com dois pontos de quadratura.

\begin{table}[h]
	\begin{center}
		\caption{Valores para $H_{q}(x)$ e $v_{k}$ para $q=2,3 e 4$}
		\label{tableHermite11cap3}\vspace*{0.7cm}
		\begin{tabular}{@{}lcccc@{}}
			\hline\hline
			q & 2 & 3 & 4  \\
			\hline\hline
			$H_{q}(x)$ & $4x^{2}-1$ & $8x^{3}-12x$  & $16x^{4}-12x^{2}+12$ \\
			$v_{k}$ & $\frac{\sqrt{\pi }}{4x^{2}}$ & $\frac{96\sqrt{\pi }}{(24x^{2}-12)^{2}}$  & $\frac{768\sqrt{\pi }}{(64x^{3}-96x)^{2}}$ \\
			\hline
		\end{tabular}
	\end{center}
\end{table}

Deve-se considerar, no entanto, que a quadratura de Gauss-Hermite tem maior efetividade no cálculo de integrais do tipo

\begin{equation}
	\label{eqnHermit123cap1}
	\int_{-\infty }^{+\infty }g(x)dx=\int_{-\infty }^{+\infty }f(x)v(x)dx
\end{equation} em que:
\begin{equation}
	\label{eqnHermit124cap1}
	v(x)=v(x;\mu, \sigma ^{2}) =\frac{1}{\sqrt{2\pi \sigma ^{2}}}\exp \left \{ -\frac{1}{2} \frac{(x-\mu )^{2}}{\sigma ^{2}}\right \}
\end{equation}

Por exemplo, ao desenhar a linha para $g(x) = exp(-(x-1)^{2})$ e posicionar os pontos e pesos para esse tipo de integral, é possível visualizar a distribuição dos pontos ao longo da região de interesse, ou seja, dentro do intervalo da função de densidade e probabilidade como mostra a figura \ref{fig21cap2}:

% \begin{figure}[h]
% 	\centering
% 	\includegraphics[width=13cm,height=5.7cm]{gauss_nao_adap_m1}
% 	\caption{(a) - \textit{(Código $\mathcal{R}$ no apêndice \ref{appendixI:gaussNaoadaptativa})}}
% 	\label{fig21cap2}
% \end{figure}

Apesar dos pontos de quadratura não estarem perfeitamente alinhados com a área de maior densidade do gráfico, ainda assim o resultado obtido ao se calcular a integral para função $g(x) = exp(-(x-1)^{2})$ está próxima do valor real para ela, como visto abaixo:

\begin{verbatim}
	> sum(quad$weights * g1(quad$nodes) * exp(quad$nodes^2))
	[1] 1.771348
	> integrate(f=g1, lower=-Inf, upper=Inf)
	1.772454 with absolute error < 1.6e-06
\end{verbatim}

Para se trabalhar sob essa perspectiva, alguns pontos devem ser observados para garantir maior acurácia nos resultados, são eles: 

\begin{itemize}
	\item Os pontos e pesos devem ser facilmente encontrados e calculados;
	\item Deve possuir uma função de densidade de probabilidade normal;
	\item O polinômio de Gauss-Hermite deve ser de baixa ordem;
	\item Os pontos devem cair na região de interesse; e
	\item Os pesos devem ser todos positivos.
\end{itemize}

Como a função $g(x)$, que dita o comportamento da curva, não exerce influência sobre a distribuição dos pontos de quadratura, é provável que, dependendo do conjunto de dados, a distribuição e posicionados dos pontos afastem-se da região de interesse da função que se deseja aproximar, como é o caso das bimodalidades e assimetrias.

Em um cenário que $g(x) = exp(-5*(x-3)^{2})$, desenhando uma assimetria na distribuição, os pontos de quadratura ficam totalmente fora da região de interesse da função, como mostra a figura \ref{fig22cap2}.

% \begin{figure}[h]
% 	\centering
% 	\includegraphics[width=13cm,height=5.7cm]{gauss_nao_adap_m2}
% 	\caption{(b) - \textit{(Código $\mathcal{R}$ no apêndice \ref{appendixJ:gaussNaoAdaptativa})}}
% 	\label{fig22cap2}
% \end{figure}

Em dados bimodais, mesmo que os pontos caiam sobre a região de interesse da função, é importante considerar o comportamento e a região de agrupamento dos dados dessa bimodalidade \ref{fig23cap2}. Em ambos os casos, ou seja, tanto na assimetria quanto na bimodalidade, o valor do cálculo da integral é afetado, produzindo resultados imprecisos.

% \begin{figure}[h]
% 	\centering
% 	\includegraphics[width=13cm,height=5.7cm]{gauss_nao_adap_m3}
% 	\caption{(c) - \textit{Código $\mathcal{R}$ no apêndice \ref{appendixK:gauss-hermite}}}
% 	\label{fig23cap2}
% \end{figure}

\begin{verbatim}
	> sum(quad$weights * g3(quad$nodes) * exp(quad$nodes^2))
	[1] 0.009721331
	> integrate(f=g3, lower=-Inf, upper=Inf)
	0.7926655 with absolute error < 0.00011
	
	> sum(quad$weights * g2(quad$nodes) * exp(quad$nodes^2))
	[1] 0.8862269
	integrate(f=g2, lower=-Inf, upper=Inf)
	0.8862269 with absolute error < 1.1e-06
\end{verbatim}

A maneira mais adequada de se tratar essas limitações é a aplicação do conceito adaptativo da quadratura de Gauss-Hermite, que seria o reescalamento e modificação dos ponto de modo que eles fiquem, em sua maioria, na região de interesse da função.

Considerando a equação \ref{eqnHermit11cap1}, no modelo adaptativo os parâmetros $\mu$ e $\sigma$, chamados aqui de $\hat{\mu } $ e $\hat{\sigma ^{2}}$, são respectivamente a média e a variância da distribuição, dados por:

\begin{equation}
	\label{eqnHermit125cap1}
	\hat{\mu }  =  \underset{x}{\arg \max g(x)} 
\end{equation} e 

\begin{equation}
	\label{eqnHermit126cap1}
	\hat{\sigma ^{2}}  =  \left [ -\frac{d^{2}}{dx^{2}} \log g(x)\right ]^{-1}\mid _{x=\hat{\mu }},
\end{equation} de modo que se for definida 

\begin{equation}
	\label{eqnHermit127cap1}
	h(x)=\frac{g(x)}{v(x;\hat{\mu }, \hat{\sigma }^{2})}.
\end{equation} 

A função pode ser reescrita de modo que 

\begin{equation}
	\label{eqnHermit128cap1}
	\int_{-\infty }^{+\infty }g(x)dx = \int_{-\infty }^{+\infty }h(x)x(x;\hat{\mu }, \hat{\sigma }^{2})dx,
\end{equation} onde $v(x;\hat{\mu }, \hat{\sigma }^{2}) = \frac{1}{\sqrt{2\pi \hat{\sigma }^{2}}}\left \{ -\frac{1}{2} \frac{(x-\hat{\mu})^{2}}{\hat{\sigma }^{2}}\right \}$, de modo a aplicação da Gauss-Hermite sobre ela \ref{eqnHermit128cap1} produz a seguinte equação: 

\begin{equation}
	\label{eqnHermit129scap1}
	\int_{-\infty }^{+\infty }g(x)dx\approx \sum_{k=1}^{q}\frac{v_{k}}{\sqrt{\pi }}h(\sqrt{2}\hat{\sigma }s_{k}+\hat{\mu })=\sqrt{2\pi \hat{\sigma }}\sum_{k=1}^{q}v_{k}^{\dotplus }g(s_{k}^{\dotplus }),
\end{equation} de modo que o comportamento dos pontos seguem a distribuição dos dados, como mostra a figura \ref{fig24cap2}

% \begin{figure}[h]
% 	\centering
% 	\includegraphics[width=9cm]{gauss_adaptativa}
% 	\caption{Pontos de quadratura sobre dados assimétrico - Gauss-Hermite adaptativa}
% 	\label{fig24cap2}
% \end{figure}

\section{OLLST com efeitos aleatórios}

A estrutura do modelo hierárquico para OLLST é dada por:

\begin{equation}
	\label{eqnEfeitosAleatorios11}
	Y_{ijk}|w_{i}\overset{iid}{\sim }OLLST(\mu _{ij},\sigma ,\lambda ,\alpha ) \textrm{ e } w_{i}\overset{iid}{\sim }N(0,\sigma ^{2}_{\beta _{0}}),
\end{equation} onde $Y_{ijk}$ é a variável resposta associada ao $i$-ésimo indivíduo $(i=1, ..., m)$ para $j=(1, ...,J)$ repetições dos tratamentos e para $k=(1,...,K)$ repetições no tempo, de modo que o $i$-ésimo indivíduo pode ser descrito pelo vetor de variáveis respostas $ Y_{i}=[Y_{i11}$ $,...,Y_{i1K};$ $Y_{i21},...,Y_{i2K},...,Y_{in_{i}1},...,Y_{in_{i}K}]$

É necessário considerar também, a possibilidade de relação entre as variáveis respostas e explicativas, ou mesmo covariáveis. Nesse contexto, é correto pressupor que todas as respostas para o mesmo individuo possam ter efeitos aleatórios comuns, denotado aqui por $w_{i}$, e que esse efeito pode ser considerado como variáveis aleatórias não observadas. Nesse caso, o modelo de regressão para dados com característica de correlacionamento pode ser expresso por: 
\begin{equation}
	\label{eqnEfeitosAleatorios12}
	Y_{ijk}=\beta _{0}+w_{i}+\underset{\textrm{efeitos longitudinais}}{\underbrace{\sum_{k=1}^{3} \omega_{k}d_{k}}} + \underset{\textrm{efeitos hierárquicos}}{\underbrace{\sum_{i=2}^{4}\sum_{k=1}^{3}\delta _{ik}d_{k}u_{i}}} + \sigma z_{ijk},
\end{equation} sendo $\delta _{ik}$ os desdobramentos dos efeitos de cada tratamento dentro de cada tempo $\omega_{i}$, apresentando efeito aleatório para o individuo $i$ de variável aleatória $z_{ijk}$ com distribuição OLLST. A relação de ausência ou presença de para valores de tempo e tratamento fica a cargo das variáveis dummies $d_{k}$ $u_{k}$ com a seguinte regra:

\begin{equation}
	\label{eqnEfeitosAleatorios13s}
	d_{k}=\left\{\begin{matrix}
		0, \textrm{se a época} \neq k\\ 
		1, \textrm{se a época} = k
	\end{matrix}\right.
\end{equation}

\begin{equation}
	u_{k}=\left\{\begin{matrix}
		0, \textrm{se o tratamento} \neq k\\ 
		1, \textrm{se o tratamento} = k
	\end{matrix}\right.
\end{equation}

Sendo $f(y_{ijk}|w_{i})$ a função de densidade condicional para $w_{i}$, de modo que $Y_{ijk}|w_{i}$ tem a função de densidade conjunta dada por $f(y_{ijk},w_{i};\theta)=f(y_{ijk}|w_{i})f(w_{i}\sigma _{\beta _{0}}^{2})$, é possível obter as seguintes estruturas:

\begin{enumerate}
	\item $Y_{ijk}|w_{i} \sim OLLST(\mu _{ij},\sigma ,\lambda ,\alpha )$ tem a função de densidade condicional expressa  por:
	\subitem \begin{equation}
		f(y_{ijk}|w_{i})=\frac{\alpha \phi (\frac{y_{ijk}-\mu _{ijk}}{\alpha })\Phi ^{\alpha -1}\left ( \frac{y_{ijk}-\mu _{ijk}}{\sigma } \right )\left [ 1-\Phi \left ( \frac{y_{ijk}-\mu _{ijk}}{\sigma } \right ) \right ]^{\alpha -1}}{\sigma \left \{ \Phi ^{\alpha } (\frac{y_{ijk}-\mu _{ijk}}{\alpha }) +\left [ 1-\Phi (\frac{y_{ijk}-\mu _{ijk}}{\alpha }) \right ]^{\alpha } \right \}^{2}}
	\end{equation}
	\item E equação \ref{eqnEfeitosAleatorios14} é a indicação de que o modelo de regressão será posto no parâmetro de forma.
	\subitem \begin{equation}
		\label{eqnEfeitosAleatorios14}
		\mu _{ijk}=\beta _{0}+w_{i}+\sum_{k=1}^{3}\omega _{k}d_{k}+\sum_{i=2}^{4}\sum_{k=1}^{3}\delta _{ik}d_{k}u_{k}
	\end{equation}
	\item A distribuição de efeitos aleatórios é definida por $w_{i}\sim N(0, \sigma _{\beta _{0}}^{2})$ e densidade:
	\subitem \begin{equation}
		\label{eqnEfeitosAleatorios15s}
		f(w_{i};\sigma _{\beta _{0}}^{2})=\frac{1}{\sigma _{\beta 0}\sqrt{2\pi }}\exp \left \{ -\frac{1}{2} \left ( \frac{w_{i}}{\sigma _{\beta _{0}}} \right )^{2}\right \},
	\end{equation} sendo $w_{i} \in \mathbb{R}\textrm{ e }\sigma _{\beta _{0}}> 0$ o intercepto de efeitos aleatórios.
\end{enumerate} 

%\section{Estudo de simulação}
%\label{cap4:estudosimulacao}
%
%Para este delineamento de simulações, tomou-se como base os trabalhos com medida repetida de \cite{park2006backfitting, guoyou2008robust, fu2012quantile} de modo que a amostra $y_{ijk}$ foi gerada de acordo com a distribuição $OLLST(\mu _{ij},\sigma ,\lambda ,\alpha )$ e a variável tempo  com valores inteiros gerados por \texttt{rep(seq(1,3, length.out = j), times = i*k)}, sendo $i=4$ e $j=3$ com variação das repetições $k$ em $k=4$, $k=8$ e $k=16$. Foram geradas 1000 amostras de tamanho $i \times j \times k$ com distribuição OLLST sob as seguintes configurações:
%
%\begin{itemize}
%	\item $\sigma = 0.5$ e $\alpha = 0.5$;
%	\item $\beta_{0}=1.5$, $\beta_{1}=0.1$ e $\beta_{2}=-0.6$;
%	\item Parâmetro de escala $\sigma_{0}=1$;
%	\item Parâmetro de forma $\nu_{0}=0,6$; e
%	\item Parâmetro de forma $\tau_{0}=0,7$.
%\end{itemize}
%
%A análise dos resíduos para o estudo de simulação de Monte Carlo tomou como base as pesquisas de  \cite{feng2017randomized, scudilio2020adjusted}, de modo que fosse possível entender do comportamento do modelo através dos ajustes das estimativas médias (EMs), vieses e erros quadráticos médios (EQMs).

\section{Envelope simulado}
%\label{cap4:envelopeSimulado}
%além de gráficos de probabilidade normal para avaliar o grau de desvio da hipótese de normalidade dos resíduos \textit{(Veja o apêndice \ref{appendixL})}.

Como parte das etapas de diagnóstico do modelo,  de modo a avaliar a veracidade da hipótese assumida para à distribuição da variável resposta \cite{fernandes2019contribuiccoes} em função do tempo, alguns gráficos de probabilidade foram gerados (envelope simulado) de banda obtidas no ajuste do modelo para diferentes conjuntos de dados sobre a variação do teor de carbono no solo.

Para construção do gráfico, considerou que o $k$-ésimo $(k=1,...,n)$ valor ordenado do resíduos versos o valor correspondente esperado da estatística de ordem normal padrão. Assim, tem-se que:

\begin{equation}
	\Phi ^{-1}\frac{k+n-\frac{1}{8}}{2n+\frac{1}{2}}
\end{equation} onde $\Phi (.)$ é a função de distribuição acumulada n $N(0,1)$.

A geração dos gráficos de probabilidade dos resíduos quantílicos foi feita através da função \texttt{halfnorm} da biblioteca \texttt{BLMM}.

\section{Aplicação: análise dos níveis de carbono no solo}
\label{cap3:secMaterial}

\subsection{Material}
\label{cap441:material}
Os dados utilizados para ajuste no modelo OLLST foram obtidos em 31 de maio de 2022 e são provenientes de um experimento realizado em Piracicaba, estado de São Paulo (SP), por \cite{gabriel2022} em sua dissertação de mestrado.

\section{Contexto dos dados}
\label{cap5:secContextoDados}

Há hoje uma preocupação genuína quanto aos modelos agrícolas de produção de alimento, principalmente ao encarar o cenário desenhado pela ONU para os próximos anos. Os modelos mais conservacionistas, como por exemplo o Sistema de Plantio Direto, que potencialmente acumula matéria orgânica no solo, tendem a se destacar frente a necessidade crescente, sobretudo pela sua característica sustentável \cite{gabriel2022}.

Estudos como \cite{kaiser2012cycling, kalbitz2000controls} ratificam essa tendência e pontuam os benefícios do acúmulo de resíduos culturais na produção de matéria orgânica do solo (MOS), bem como de matéria orgânica dissolvida (MOD ou COD), e seus efeitos na produção de carbono orgânico dissolvido (COD). Há ainda uma relação direta no transporte de nutrientes para microrganismos que, por consequência, está ligada a dinâmica de nutrientes do solo, tendo inclusive impactos a longo prazo, quando comparadas com fatores ambientais e antropogênicos. 

Considerando a lacuna existente em pesquisas voltadas para solos tropicais, \cite{gabriel2022} propôs um experimento a fim de entender o fluxo de COD e a dinâmica dele no perfil de solo Brasileiro com o objetivo de incitar adequações nas práticas de manejo adequado do solo, mitigando perda de nutrientes e possíveis contaminações de corpos d'água, com olhar voltado aos sistemas agrícolas.

Com intuito de refinar o escopo desta pesquisa, uma análise exploratória foi realizada sob o conjunto de dados a fim de identificar aquele com maior dispersão na organização das unidades amostrais. As figuras \ref{fig7differentGraphs4} e \ref{fig7differentGraphs5} mostram os boxplots para os conjunto de dados utilizados por \cite{gabriel2022} que contém a variável "tempo/época".

Para que seja viável a análise de dados em um modelo hierárquico de uma variável resposta em função do tempo, tem-se a pressuposição de independência entre as observações para o mesmo indivíduo e dependência das unidades amostrais retiradas das unidades experimentais \cite{cnaan1997using, fausto2008modelo}, necessitando, por tanto, uma avaliação primária de cada conjunto.

\subsection{Descrição dos Dados}
\label{cap3:secDescricaodosdaos}

Nesta pesquisa foram ajustados modelos de regressão para seis conjunto diferentes de dados sobre a dinâmica do carbono em camadas de Latossolo Vermelho (VL). Apesar de focar ná análise do carbono, \cite{gabriel2022} separou alguns conjuntos de dados para a fim de entender as influências que o carbono poderia sofre durante o experimento. Sete foram os conjuntos de dados, no entanto, um não dispunha da variável tempo. As nomenclaturas para cada conjunto seguirão como:
 
 \begin{itemize}
 	\item \textbf{LV-C} - conjunto de dados para análise do teor de carbono em função do tempo e profundidade; 
 	\item \textbf{LV-CN} - conjunto de dados para análise da razão entre carbono e nitrogênio; 
 	\item \textbf{LV-COD} - conjunto de dados para análise do fluxo de carbono dissolvido; 
 	\item \textbf{LV-EC} - conjunto de dados para análise do teor de carbono em função do tempo; 
 	\item \textbf{LV-EN} - conjunto de dados para análise do teor de carbono em função do tempo; e
 	\item \textbf{LV-N} - conjunto de dados para avaliar níveis de nitrogênio.
 \end{itemize}
 

Cada conjunto de dados tem as seguintes variáveis:

\begin{itemize}
	\item  \textbf{Trat.} Variável categórica referente ao tipo de tratamento aplicado: \textbf{(M)} Palha de milho (Zea mays); \textbf{(S)} Palha soja (Glycine max); \textbf{(MS)} A combinação das palhas de soja e milho; \textbf{(SR)} O controle sem nenhum tipo de tratamento; 
	\item \textbf{Epoca}[\textit{sic}]. Variável categórica referente ao tempo de observação da unidade amostral. ($t$1, $t$2, $t$3); e
	\item \textbf{Y.} Variável resposta referente ao teor de carbono no solo para aquela unidade experimental.
\end{itemize}

O conjunto \textbf{LV-C} possui uma população de tamanho 144 sendo: 3 variações de época $t=(1,2,3)$, 4 tipos de tratamentos $k=(1,2,3,4)$ e 3 profundidades $x=(010, 020, 030)$; o conjunto \textbf{LV-CN} possui uma população de tamanho 144 sendo: 3 variações de época $t=(1,2,3)$, 4 tipos de tratamentos $k=(1,2,3,4)$ e 3 profundidades $x=(010, 020, 030)$; o conjunto \textbf{LV-COD} possui uma população de tamanho 176 sendo: 11 variações de época $t=(1,2,3,4,5,6,7,8,9,10,11)$ e 4 tipos de tratamentos $k=(1,2,3,4)$; o conjunto \textbf{LV-EC} possui uma população de tamanho 48 sendo: 3 variações de época $t=(1,2,3)$ e 4 tipos de tratamentos $k=(1,2,3,4)$; o conjunto \textbf{LV-EN} possui uma população de tamanho 48 sendo: 3 variações de época $t=(1,2,3)$ e 4 tipos de tratamentos $k=(1,2,3,4)$; e o conjunto \textbf{LV-N} possui uma população de tamanho 144 sendo: 3 variações de época $t=(1,2,3)$, 4 tipos de tratamentos $k=(1,2,3,4)$ e 3 profundidades $x=(010, 020, 030)$. Para todos os modelos foram realizadas 4 repetições das observações.

Vale ressaltar que para seu estudo, \cite{gabriel2022} considerou também a profundidade (em cm) para as amostras retiradas. No entanto, como o foco deste trabalho reside na análise de dados com medida repetida no tempo, está variável não foi incluída nos experimentos.

As figuras \ref{fig7differentGraphs1}, \ref{fig7differentGraphs3} e \ref{fig7differentGraphs2} mostram a distribuição da frequência do carbono orgânico analisado por \cite{gabriel2022}. Cada gráfico trás consigo uma linha pontilhada em azul que representa uma média para variável resposta, além de um mancha em vermelho para representar a zona de densidade mostrando a variação nas distribuições para cada conjunto.

% \begin{figure}[!htb]
% 	\begin{center}
% 		\hspace{0.2cm}(LV-C)\hspace{5cm} (LV-CN)\\
% 		\includegraphics[height=4.7cm]{LV_C_HIST}
% 		~~\includegraphics[height=4.7cm]{LV_CN_HIST}\\
% 		\caption{Concentração de carbono no solo.}
% 		\label{fig7differentGraphs1}
% 	\end{center}
% \end{figure}

% \begin{figure}[!htb]
% 	\begin{center}
% 		\hspace{0.2cm}(LV-EN)\hspace{5cm} (LV-EC)\\
% 		\includegraphics[height=4.7cm]{LV_EN_HIST}
% 		~~\includegraphics[height=4.7cm]{LV_EC_HIST}\\
% 		\caption{Concentração de carbono no solo.}
% 		\label{fig7differentGraphs3}
% 	\end{center}
% \end{figure}

% \begin{figure}[!htb]
% 	\begin{center}
% 		\hspace{0.2cm}(LV-COD)\hspace{5cm} (LV-N)\\
% 		\includegraphics[height=4.3cm]{LV_COD_HIST}
% 		~~\includegraphics[height=4.3cm]{LV_N_HIST}\\
% 		\caption{Concentração de carbono no solo.}
% 		\label{fig7differentGraphs2}
% 	\end{center}
% \end{figure}

\section{Análise Exploratória}
\label{cap5:secAnaliseExpliratoria}

Considerando os valores para cada tipo de controle em razão da época, verificou-se que conjunto \textbf{LC-V}, conforme a figura \ref{cap3:figLVC}, sugerem alterações nos valores para variável resposta em razão do tempo, justificando uma avaliação do modelo para cada tipo de controle (tratamento) dentro de suas respectivas épocas;

% \begin{figure}[h]
% 	\centering
% 	\includegraphics[width=13cm]{BOX_LV_C}
% 	\caption{Boxplot para o conjunto LV-C}
% 	\label{cap3:figLVC}
% \end{figure} 

% \begin{figure}[h]
% 	\centering
% 	\includegraphics[width=13cm]{BOX_LV_CN}
% 	\caption{Boxplot para o conjunto LV-CN}
% 	\label{cap3:figLVCN}
% \end{figure} 

Para o conjunto \textbf{LV-CN} o boxplot mostrado na figura \ref{cap3:figLVCN} exibe a presença de valores extremantes para alguns tratamentos no segundo e terceiro momento da observação. \textit{Outliers} também são observados nos conjuntos de dados \textbf{LV-COD} (figura \ref{cap3:figLVCOD}), \textbf{LV-EC} (figura \ref{cap3:figLVEC}) e \textbf{LV-EN} (figura \ref{cap3:figLVEN}). Importante ressaltar que, segundo \cite{barnett1994outliers}, o \textit{outlier} é uma observação (ou conjunto delas) que aparenta ser inconsistente com o restante do conjunto de dados, com crucial importância para caracterização da amostra, uma vez que, na ausência deles o modelo pode apresentar um comportamento diferente. Em todos os casos, a interpretação sugere que, ainda que mínima, o tempo exerce influência no comportamento da variável resposta, justificando uma análise em segundo nível.  

% \begin{figure}[h]
% 	\centering
% 	\includegraphics[width=13cm]{BOX_LV_COD}
% 	\caption{Boxplot para o conjunto LV-COD}
% 	\label{cap3:figLVCOD}
% \end{figure} 


% \begin{figure}[h]
% 	\centering
% 	\includegraphics[width=13cm]{BOX_LV_EC}
% 	\caption{Boxplot para o conjunto LV-EC}
% 	\label{cap3:figLVEC}
% \end{figure} 

% \begin{figure}[h]
% 	\centering
% 	\includegraphics[width=13cm]{BOX_LV_EN}
% 	\caption{Boxplot para o conjunto LV-EN}
% 	\label{cap3:figLVEN}
% \end{figure}

Assim como o conjunto \textbf{LV-C}, o \textbf{LV-N} mostrou relativo controle entre os valores da variável resposta, sem a presença de pontos extremos e tendência que sugerem variação no valor para os tipos de controle em razão do tempo, de modo a valer uma análise mais aprofundada da interação entre as variáveis.
 
% \begin{figure}[h]
% 	\centering
% 	\includegraphics[width=13cm]{BOX_LV_N}
% 	\caption{Boxplot para o conjunto LV-N}
% 	\label{cap3:figLVN}
% \end{figure} 

\section{Tratamento de saída da biblioteca GAMLSS}
\label{cap3:secTratamentoSaidaGamlss}

Durante análise do modelo OLLST, incorporado a família de distribuições do GAMLSS, verificou-se uma característica da saída de resultados para o modelo ajustado que poderia representar uma dificuldade na interpretação. A implementação computacional do modelo descrito na equação \ref{eqnEfeitosAleatorios12} implica em uma saída hierarquizada onde os tratamentos são comparados dentro de cada dimensão (tempo).

Supondo que se pretende ajustar um modelo de regressão para dados com medida repetida no tempo, cuja a estrutura da tabela seja semelhante ao mostrado na tabela \ref{tab1cap2disposicao}. 
Ao se descrever computacionalmente o modelo da equação \ref{eqnEfeitosAleatorios12}, perceber-se que cada comparação entre indivíduos/tratamentos fica subordinada a variável \textbf{tempo} \textit{(Vide a figura \ref{fig24cap21})}.

% \begin{figure}[h]
% 	\centering
% 	\includegraphics[width=15cm]{tempo vs tratamentos}
% 	\caption{Representação do modelo em função do tempo}
% 	\label{fig24cap21}
% \end{figure}

Outro fator que deve ser considerado neste cenário está relacionado a indicação daquilo que deve ser comparado. Em uma estrutura com $k$ indivíduos/tratamentos, cada qual deve passar pela mesma rotina de comparação, ou seja, o dispositivo apresentado na figura \ref{fig24cap21} será repetido em $\sum_{i=1}^{k}$, e cada interação terá sua própria saída de dados. 

Na postulação do modelo hierárquico a identificação do indivíduo/tratamento em restrição, ou seja, aquele que deverá ser comparado com os demais para cada lâmina do tempo, não fica evidenciado na saída dos resultados. Nesta ocasião, o pesquisador deve inferir qual indivíduo/tratamento está em restrição, como mostra a figura \ref{fig24cap22}, para somente então poder vislumbrar o panorama comparativo completo. 

% \begin{figure}[h]
% 	\centering
% 	\includegraphics[width=16cm]{saida}
% 	\caption{Saída do ajuste para determinado tratamento}
% 	\label{fig24cap22}
% \end{figure} 

Ademais, a indicação manual para cada restrição pode trazer uma camada extra de dificuldade para o pesquisador, caso o conjunto de dados tenha muitos tratamentos. Consequentemente, a automação desta rotina representaria uma significativa economia de tempo, além de minimizar a possibilidade de erros. 

Por fim, deve-se considerar a pluralidade da estatística e suas aplicações em diversos contextos científicos. Desenvolvida com propósitos de manipulação e análise de dados \cite{martins2016programaccao} a linguagem de programação $\mathcal{R}$ pode parecer desafiadora para quem não teve contato com ela, mesmo para aqueles que já se aventuraram no munda da programação. 

Em $\mathcal{R}$, grande parte da manipulação dos dados é feita pela codificação de instruções, fazendo-se poucas interações com algum tipo de interface gráfica. Mesmo não sendo um impeditivo para o uso da linguagem, a interação gráfica de rotinas já modeladas pode ampliar seu uso tornando, por exemplo, o ajuste de modelos mais intuitivo. 

Pensando nisso, parte dos esforços deste trabalho concentraram-se na proposta de uma melhoria para saída de resultados de modelos hierárquicos pertencentes à família de distribuição GAMLLS, com a possibilidade ou não de manipulação gráfica.

As tecnologias utilizadas para proposta foram:

\begin{itemize}
	\item Linguagem de programação $\mathcal{R}$ para lapidação da saída do modelo ajustado, tendo como requisitos as seguintes bibliotecas:
		\subitem \textbf{nlme} - para trabalho com modelos de efeito linear misto;
		\subitem \textbf{hnp} - para análise da qualidade do ajuste do modelo;
		\subitem \textbf{gamlss} - para ajuste de modelo;
		\subitem \textbf{dplyr} - para manipulação de dados;
		\subitem \textbf{tidyr} - para manipulação e organização de dados;
		\subitem \textbf{ggplot2} - para plotagem de gráficos;
		\subitem \textbf{string} - para manipulação de strings; e
		\subitem \textbf{rjson} - manipulação de arquivos \texttt{.json};
	\item Linguagem de programação Python, para criação das rotinas de backend como orquestração de scripts externos e captura de parâmetros. Para isso foi necessário a instalação das seguinte bibliotecas:
		\subitem \textbf{Flask==2.3.2} - micro framework para programação web com python; e
		\subitem \textbf{Jinja2==3.1.2} - manipulação de instruções python dentro de arquivos HTML.
	
	\item Shell scritp - para automação das rotinas de variáveis de ambiente do sistema operacional Linux;
	\item HTML , CSS e JS - para construção da interface de interação do usuário com modelo; e
	\item Json - para armazenamento temporário dos parâmetros de chamada do scritp r.
\end{itemize}


\section{Observações finais}
\label{cap3:observacoesFinais}

Este capítulo apresentou o contexto de criação do modelo de regressão OLLST, incorporado a família de distribuição GAMLSS, com propósito de trazer uma maior flexibilidade nos ajustes de parâmetros, concentrando-se não apenas na média, mas aceitando ajuste para todos os parâmetros da função. Como uma extensão da Skew $t$-Student, através da variação dos valores $\alpha$, $\lambda$ e $\nu$, o modelo possibilita o ajuste para uma maior amplitude de comportamentos \cite{dos2021overview}, graças a sua função de estimação pelo método da máxima verossimilhança. 

Considerando o cenário favorável proposto por \cite{adolf2021}, mostrou-se aqui também a descrição matemática do modelo OLLST para dados longitudinais, de modo que seja possível analisá-lo com dados de medida repetida.

Ao final, foi discorrido sobre uma proposta de lapidação para saída de dados dos modelos hierárquicos ajustado pela função \texttt{gamlss} do $\mathcal{R}$. Com o objetivo de tornar mais intuitivo e autônomo o processo execução e análise dos dados, foi sugerido uma ferramenta capaz de automatizar as rotinas de comparação para diferentes conjuntos de dados onde o usuário, por meio de uma interface gráfica, seja capaz de informar o conjunto de dados, colunas que deseja utilizar e quais os contextos de variáveis que elas representam, informando por fim a distribuição de dados para qual deseja ajustar o modelo. Para tal, foram descritas as ferramentas e tecnologias utilizados no processo de criação desta ferramenta.

\chapter{Resultados e Discussões}
\label{cap4:title}

\section{Análise descritiva}
%A análise do modelo hierárquico tomou como base a equação 
%\begin{equation}
%	\label{eqnEfeitosAleatorios41}
%	Y_{ijk}=\beta _{0}+w_{i}+\underset{\textrm{efeitos longitudinais}}{\underbrace{\sum_{k=1}^{3} w_{k}d_{k}}} + \underset{\textrm{efeitos hierárquicos}}{\underbrace{\sum_{i=2}^{4}\sum_{k=1}^{3}\delta _{ik}d_{k}u_{i}}} + \sigma z_{ijk}
%\end{equation} sendo $\delta _{ik}$ os desdobramentos dos efeitos de cada tratamento dentro de cada tempo $\omega_{i}$, $i$ o vetor de indivíduos, $j$ o vetor de repetição para cada individuo e $k$ o vetor de medida repetida no tempo para cada tratamento. Além disso, o modelo conta com a variável de efeito aleatório para cada individuo, $z_{ijk}$, com distribuição OLLST. 

Para o conjunto de dados \textbf{LV-C}, cujo o gráfico de frequência é apresentado na figura \ref{fig7differentGraphs1}(LV-C), utilizando a função \texttt{histDist} é possível visualizar uma perspectiva de concentração para zona de densidade de probabilidade nas famílias de distribuição Normal e OLLST, como mostra a figura \ref{fig7differentGraphs8}. Através da saída para esta função, apresentada aqui na tabela \ref{tab8cap5histDist} é possível inferir, preliminarmente, um melhor desempenho da distribuição OLLST em detrimento a NO para esse conjunto de dados.

% \begin{figure}[!htb]
% 	\begin{center}
% 		\hspace{0.2cm}(NO)\hspace{5cm} (OLLST)\\
% 		\includegraphics[height=4.7cm]{histdistno}
% 		~~\includegraphics[height=4.7cm]{histdistollst}\\
% 		\caption{Conjunto de dados LV-C}
% 		\label{fig7differentGraphs8}
% 	\end{center}
% \end{figure}

\begin{table}[h]
	\begin{center}
		\caption{Saída para função \texttt{histDist()} - Conjunto LV-C}\vspace*{0.3cm}
		\label{tab8cap5histDist}\vspace*{0.3cm}
		\begin{tabular}{@{}lrr@{}}
			\toprule
			& \multicolumn{2}{c}{Distribuição} \\ \cmidrule(l){2-3} 
			& Normal          & OLLST           \\ \midrule
			$\mu$               & 31.96           & 24.96           \\
			$\sigma$           & 2.059            & 1.928            \\
			$\nu$           & -            & 4.883           \\
			$\tau$           & -            & -0.7497            \\
			Global Deviance & 1001.56         & 948.102         \\
			AIC             & 1005.56         & 956.102         \\
			SBC             & 1011.5         & 967.981         \\ \bottomrule
		\end{tabular}
	\end{center}
\end{table}

Submetido ao ajuste pelo modelo OLLST, o conjunto de dados \textbf{LV-C} teve como saída os resultados apresentados na tabela.As hipóteses formuladas para o teste de comparação levou em consideração cada época (1,2 e 3) a fim de entender se havia ou não diferença significativa entre os tipos de controle. Para o nível de significância de 5\% tem-se que:
\begin{itemize}
	\item Na época 1 não se constatou uma diferença significativa entre os tratamento;
	\item Na época 2, o tratamento 3 quando comparado ao tratamento 4 não apresentou diferença significativa ($p$-valor = 0.8387);
	\item A terceira época mostra que há diferenças significativas nas comparações entre todos os tratamentos; e
	\item Os melhores resultados obtidos, em função do tempo, encontram-se na época 3.
\end{itemize} \ref{tableResultados11}. 

\begin{table}[h]
	\begin{center}
		\caption{Comparação dos efeitos de cada tratamentos}
		\label{tableResultados11}\vspace*{0.3cm}
		\begin{tabular}{@{}lccccc@{}}
			\hline\hline
			Época & Hipóteses & Estimativa & Std. Error &  $p$-valor  \\
			\hline\hline
			1 & H$_{0}: \tau_{2} - \tau_{1}=0$ 	& 0.000000$^{ns} $ 		& 0.3113539 	&  1.000 \\
			1 & H$_{0}: \tau_{3} - \tau_{1}=0$ 	& 0.000000$^{ns} $ 		& 0.3113539 	&  1.000 \\
			1 & H$_{0}: \tau_{4} - \tau_{1}=0$ 	& 0.000000$^{ns} $ 		& 0.3113539 	&  1.000 \\
			1 & H$_{0}: \tau_{3} - \tau_{2}=0$ 	& 0.000000$^{ns} $ 		& 0.3113538 	&  1.000 \\
			1 & H$_{0}: \tau_{4} - \tau_{2}=0$ 	& 0.000000$^{ns} $ 		& 0.3113538 	&  1.000 \\
			1 & H$_{0}: \tau_{4} - \tau_{3}=0$ 	& 0.000000$^{ns} $ 		& 0.3113539 	&  1.000 \\
			\hline
			2 & H$_{0}: \tau_{2} - \tau_{1}=0$ 	& 2.224100$^{*} $ 		& 0.5784348 	&  0.0002\\
			2 & H$_{0}: \tau_{3} - \tau_{1}=0$ 	& -2.585490$^{*} $ 		& 0.4519209 	&  0.0000\\
			2 & H$_{0}: \tau_{4} - \tau_{1}=0$ 	& -2.485135$^{*} $ 		& 0.3846374 	&  0.0000\\
			2 & H$_{0}: \tau_{3} - \tau_{2}=0$ 	& -4.809589$^{*} $ 		& 0.6547566 	&  0.0000\\
			2 & H$_{0}: \tau_{4} - \tau_{2}=0$ 	& -4.709235$^{*} $ 		& 0.6102620 	&  0.0000\\
			2 & H$_{0}: \tau_{4} - \tau_{3}=0$ 	& 0.100354$^{ns} $ 		& 0.4920016		&  0.8387 \\
			\hline
			3 & H$_{0}: \tau_{2} - \tau_{1}=0$ 	& 7.064183$^{*} $ 		& 0.2837398 	&  0.0000\\
			3 & H$_{0}: \tau_{3} - \tau_{1}=0$ 	& 7.791710$^{*} $ 		& 0.2705048 	&  0.0000\\
			3 & H$_{0}: \tau_{4} - \tau_{1}=0$ 	& 10.500265$^{*}$ 		& 0.3903520 	&  0.0000\\
			3 & H$_{0}: \tau_{3} - \tau_{2}=0$ 	& 0.727528$^{*} $ 		& 0.2011386 	&  0.0004\\
			3 & H$_{0}: \tau_{4} - \tau_{2}=0$ 	& 3.436082$^{*} $ 		& 0.3459170 	&  0.0000\\
			3 & H$_{0}: \tau_{4} - \tau_{3}=0$ 	& 2.708554$^{*} $ 		& 0.3351465 	&  0.0000\\
			\hline\hline \\
			$^{1} $ & \multicolumn{4}{c}{O rótulo $^{*}$ indica que há diferença significava entre os tratamentos considerando;} \\
			$^{1} $ & \multicolumn{4}{c}{O rótulo $^{ns}$ indica que não há diferença significava entre os tratamentos.} \\ 
		\end{tabular}
	\end{center}
\end{table}

Na avaliação do comportamento do modelo OLLST em comparação ao NO, através do gráfico de probabilidade mostrados na figura \ref{fig7differentGraphs15} é possível concluir para esse conjunto de dados que o modelo OLLST saiu-se melhor no ajunte, com 25\% dos pontos fora da banda enquanto o NO teve 29,17\%.

% \begin{figure}[!htb]
% 	\begin{center}
% 		\hspace{0.2cm}(NO)\hspace{5cm} (OLLST)\\
% 		\includegraphics[height=4.5cm]{resido_NO}
% 		~~\includegraphics[height=4.5cm]{residos_OLLST}\\
% 		\caption{Conjunto de dados LV-C}
% 		\label{fig7differentGraphs15}
% 	\end{center}
% \end{figure}

Ainda que o OLLST tenha se saído melhor na comparação acima, esta afirmação restringe-se ao contexto do conjunto de dados LV-C. Por tanto, uma nova análise foi realizada desta vez com o conjunto \textbf{LV-CN}, mostrado na figura \ref{fig7differentGraphs1}(LV-CN) seu gráfico de frequência. 

Novamente as informações sobre o conjunto obtidas através da função \texttt{histDist} (veja a figura \ref{fig7differentGraphs16}) levam a crer que a distribuição OLLST encaixa-se melhor a zona de densidade de probabilidade da população, informações que corroboram com a saída de dado mostradas na tabela \ref{tab8cap5histDist1}. No entanto, o gráfico de probabilidade para o conjunto mostra um ajuste mais preciso para distribuição NO, com zero pontos fora das bandas do envelope simulado. 

% \begin{figure}[!htb]
% 	\begin{center}
% 		\hspace{0.2cm}(NO)\hspace{5cm} (OLLST)\\
% 		\includegraphics[height=4.5cm]{LV_CN_HIST_NO}
% 		~~\includegraphics[height=4.5cm]{LV_CN_HIST_OLLSTT}\\
% 		\caption{HistDist para conjunto LV-CN}
% 		\label{fig7differentGraphs16}
% 	\end{center}
% \end{figure}

\begin{table}[h]
	\begin{center}
		\caption{Saída para função \texttt{histDist()} - Conjunto LV-CN}\vspace*{0.3cm}
		\label{tab8cap5histDist1}\vspace*{0.3cm}
		\begin{tabular}{@{}lrr@{}}
			\toprule
			& \multicolumn{2}{c}{Distribuição} \\ \cmidrule(l){2-3} 
							& Normal        & OLLST           \\ \midrule
			$\mu$           & 11.38         & 8.007           \\
			$\sigma$        & 0.05257       & 4.8            \\
			$\nu$           & -            	& 0.03472           \\
			$\tau$          & -            	& 4.907            \\
			Global Deviance & 423.795       & 414.762         \\
			AIC             & 427.795       & 422.762         \\
			SBC             & 433.735       & 434.641         \\ \bottomrule
		\end{tabular}
	\end{center}
\end{table}

% \begin{figure}[!htb]
% 	\begin{center}
% 		\hspace{0.2cm}(NO)\hspace{5cm} (OLLST)\\
% 		\includegraphics[height=4.5cm]{LV_CN_ENV_NO}
% 		~~\includegraphics[height=4.5cm]{LV_CN_ENV_OLLST}\\
% 		\caption{Envelope simulado para o conjunto LV-CN}
% 		\label{fig7differentGraphs14}
% 	\end{center}
% \end{figure}


Deferente dos demais conjuntos, onde os tempos de observação foram somente três, para o conjunto \textbf{LV-COD} as amostras foram observadas em onze momentos diferentes. Considerando que estes dados avaliam a dissolução do carbono no solo, o gráfico \ref{fig7differentGraphs2}(LV-COD) mostra um comportamento em decréscimo dos valores para variável resposta com uma zona de concentração assimétrica. 

Nesse conjunto, a área sob a curva para distribuição normal (veja a figura \ref{fig7differentGraphs17}(NO)), mesmo tendendo para zona de densidade de probabilidade da população, não conseguiu abranger uma maior parcela, enquanto a OLLST (veja a figura \ref{fig7differentGraphs17}(OLLST)) elevou sua moda mais próxima ao pico da zona. Como resultado, cada distribuição apresentou valores mostrados na tabela \ref{tab8cap5histDist4}.


% \begin{figure}[!htb]
% 	\begin{center}
% 		\hspace{0.2cm}(NO)\hspace{5cm} (OLLST)\\
% 		\includegraphics[height=4.2cm]{LV_COD_NO_HISTDIST}
% 		~~\includegraphics[height=4.2cm]{LV_COD_OLLST_HISTDIST}\\
% 		\caption{HistDist para o conjunto LV-COD}
% 		\label{fig7differentGraphs17}
% 	\end{center}
% \end{figure}

\begin{table}[h]
	\begin{center}
		\caption{Saída para função \texttt{histDist()} - Conjunto LV-COD}\vspace*{0.3cm}
		\label{tab8cap5histDist4}\vspace*{0.3cm}
		\begin{tabular}{@{}lrr@{}}
			\toprule
			& \multicolumn{2}{c}{Distribuição} \\ \cmidrule(l){2-3} 
			& Normal        & OLLST           \\ \midrule
			$\mu$           & 3.087         	& 1.305           \\
			$\sigma$        & 1.167       	& 1.137            \\
			$\nu$           & -            	& 9.391           \\
			$\tau$          & -            	& -0.6819            \\
			Global Deviance & 910.21       & 737.185         \\
			AIC             & 914.21       & 745.185         \\
			SBC             & 920.551       & 757.867         \\ \bottomrule
		\end{tabular}
	\end{center}
\end{table}

O indicativo de melhora na performance do OLLST em relação ao NO na análise acima foi contatado no gráfico de probabilidade como mostra a figura \ref{fig7differentGraphs10}. No entanto, ambos os modelos apresentaram um desempenho aquém com mais da metade dos pontos fora dos limites do envelope simulado. Dos 176 pontos, 134 escaparam do intervalo das bandas para o modelo NO, enquanto OLLST foram 104. 


% \begin{figure}[!htb]
% 	\begin{center}
% 		\hspace{0.2cm}(NO)\hspace{5cm} (OLLST)\\
% 		\includegraphics[height=4.2cm]{envelope_lv_cod_no}
% 		~~\includegraphics[height=4.2cm]{envelope_lv_cod_ollst}\\
% 		\caption{Envelope simulado para o conjunto LV-COD}
% 		\label{fig7differentGraphs10}
% 	\end{center}
% \end{figure}


Na análise do conjunto de dados \textbf{LV-EN}, cujo o histograma é dado pela figura \ref{fig7differentGraphs3}(LV-EN), apesar de uma tendência da moda para região de pico da gráfico na análise OLLST (veja a figura \ref{fig7differentGraphs11}(OLLST)), ao ajustar o modelo de regressão hierárquico verificou-se que o Normal saiu-se melhor, apresentando um ajuste mais assertivo, como mostra a figura \ref{fig7differentGraphs12}(NO) e \ref{fig7differentGraphs12}(OLLST). Este comportamento pode ser reflexo da variação abrupta nos valores para variável resposta de uma população pequena, já que esse conjunto é uma parcela do estudo geral proposto por \cite{gabriel2022} para analisar o teor de carbono no solo em função do tempo de uma única profundidade. 

% \begin{figure}[!htb]
% 	\begin{center}
% 		\hspace{0.2cm}(NO)\hspace{5cm} (OLLST)\\
% 		\includegraphics[height=4.5cm]{LV_EN_HISTDIST_NO}
% 		~~\includegraphics[height=4.5cm]{LV_EN_HISTDIST_OLLST}\\
% 		\caption{HistDist para o conjunto LV-EN}
% 		\label{fig7differentGraphs20}
% 	\end{center}
% \end{figure}

\begin{table}[h]
	\begin{center}
		\caption{Saída para função \texttt{histDist()} - Conjunto LV-EN}\vspace*{0.3cm}
		\label{tab8cap5histDist6}\vspace*{0.3cm}
		\begin{tabular}{@{}lrr@{}}
			\toprule
			& \multicolumn{2}{c}{Distribuição} \\ \cmidrule(l){2-3} 
			& Normal        & OLLST           \\ \midrule
			$\mu$           & 6.185         & 5.459           \\
			$\sigma$        & -0.192       	& -0.05951            \\
			$\nu$           & -            	& 2.504           \\
			$\tau$          & -            	& -0.3176             \\
			Global Deviance & 117.783       & 113.113        \\
			AIC             & 121.783       & 121.113         \\
			SBC             & 125.525       & 128.598          \\ \bottomrule
		\end{tabular}
	\end{center}
\end{table}

% \begin{figure}[!htb]
% 	\begin{center}
% 		\hspace{0.2cm}(NO)\hspace{5cm} (OLLST)\\
% 		\includegraphics[height=5.5cm]{envelope_lv_en_no}
% 		~~\includegraphics[height=5.5cm]{envelope_lv_en_ollst}\\
% 		\caption{Envelope simulado para o conjunto LV-EN}
% 		\label{fig7differentGraphs12}
% 	\end{center}
% \end{figure}

Comportamento semelhante ao observado no ajuste para os dados \textbf{LV-EN} ocorreu com \textbf{LV-EC}. A figura \ref{fig7differentGraphs3}(LV-EC) mostra a mesma tendência dos dados, em uma proporção menor, onde há uma varição acentuada nos valores dos níveis de carbono em um conjunto com apenas quatro tratamentos, três tempos e quatro repetições. Como resultado, novamente o modelo NO ajusta-se melhor o comportamento dos dados como mostra o gráfico \ref{fig7differentGraphs11}(NO) e \ref{fig7differentGraphs11}(OLLST), onde 46 dos 48 pontos permaneceram dentro de intervalo das bandas para o modelo NO, ao passo que no OLLST esse número subiu para 27.


% \begin{figure}[!htb]
% 	\begin{center}
% 		\hspace{0.2cm}(NO)\hspace{5cm} (OLLST)\\
% 		\includegraphics[height=5.5cm]{envelope_lv_ec_no}
% 		~~\includegraphics[height=5.5cm]{envelope_lv_ec_ollst}\\
% 		\caption{Conjunto de dados LV-EC}
% 		\label{fig7differentGraphs11}
% 	\end{center}
% \end{figure}

Finalmente, a mesma sequência de experimentos foi executada para o conjunto de dado \textbf{LV-N}. Conforme visto na figura \ref{fig7differentGraphs2}(LV-N), o conjunto parece apresentar uma espécie de bimodalidade, onde picos das modas não apresentam uma acentuada disparidade nos valores para variável resposta. Como resultado dos testes o modelo OLLST saiu-se melhor quando comparado ao modelo NO, como mostra a saída para função \texttt{summary} para ambos os modelos, assim como seus respectivos gráficos de probabilidade (veja a figura \ref{fig7differentGraphs13}). No envelope simulado, apenas 7,64\% dos 144 pontos não ajustaram-se bem pelo modelo OLLST, enquanto para o NO essa porcentagem chegou 52,08\% do total.

\begin{table}[h]
	\begin{center}
		\caption{Summary of Comorbidities.}
		\label{tab10cap5summaryAllByComorbidities}\vspace*{0.3cm}
		\begin{tabular}{@{}lrrrrrrr@{}}
			\toprule
			& \multicolumn{3}{c}{Normal} & \multicolumn{1}{l}{} & \multicolumn{3}{c}{OLLST} \\ \cmidrule(lr){2-4} \cmidrule(l){6-8} 
			& $\mu$        &   & $\sigma$  &                      & $\mu$        &   & $\sigma$ \\ \midrule
			Pr(>|t|)       	& <2e-16$^{***}$   	&  & 1.8e-11$^{***}$  &       & <2e-16$^{***}$ &   	& 9.95e-05$^{***}$  \\
			Std.            & 0.05415     		&	& 0.05893   		&       & 0.05728     	&		& 0.1303  \\
			Estimate         & 2.82747       	&	& -0.43118  		&       & 2.20965    	&		& -0.5224    \\
			Global Deviance & 284.4751          & 	&        		&       & 236.3024  		&	&       \\
			AIC             & 292.5226          & 	&        		&       & 249.7433  		&	&       \\
			SBC             & 304.4724          & 	&        		&       & 269.7018  		&	&       \\ \midrule
			\multicolumn{8}{l}{\begin{tabular}[c]{@{}l@{}}Signif. codes: \\
					0 `***' 0.001 `**' 0.1 `*' 0.05 `.' 0.1 ` ' 1; $ns$ - not significant\end{tabular}}    \\ \bottomrule
		\end{tabular}
	\end{center}
\end{table} 

% \begin{figure}[!htb]
% 	\begin{center}
% 		\hspace{0.2cm}(NO)\hspace{5cm} (OLLST)\\
% 		\includegraphics[height=5.5cm]{envelope_lv_n_no}
% 		~~\includegraphics[height=5.5cm]{envelope_lv_n_ollst}\\
% 		\caption{Envelope simulado para o conjunto LV-N}
% 		\label{fig7differentGraphs13}
% 	\end{center}
% \end{figure}



\section{Proposta de melhoria para saída do modelo hierárquico} 

\subsection{Automação e melhora da saída}
\label{cap3:secAutomacaoMelhora}

A primeira etapa do processo proposto na seção \ref{cap3:secTratamentoSaidaGamlss} foi a construção de um \textit{scritp} $\mathcal{R}$ capaz de se adaptar a variações de conjunto de dados. Como parâmetro de execução, o \textit{scritp} lê um arquivo json \textit{(vide o apêndice \ref{appendixD})} onde são informados as colunas do dataframe que serão utilizadas, assim como o nome do arquivo e a distribuição da família GAMLSS que deverá ser aplicada para o ajuste do modelo.

O arquivo também é capaz de identificar as interações possíveis para cada indivíduo, criando suas próprias restrições de forma automática, gerando uma saída de dados em \texttt{.csv}, além de um gráfico de perfis da variável resposta em função do tempo. Contudo, a maior parte das funções descrita no arquivo são destinadas ao tratamento da saída do ajuste, de modo a facilitar a interpretação por parte do pesquisador. Deste modo, cada restrição foi tratada de forma a explicitar os indivíduos/tratamentos da comparação em cada tempo, mostrando ainda se esta comparação foi significativa ou não. 

\subsection{Manipulação gráfica}
\label{cap3:secManipulacaoGrafica}

Como parte da melhoria proposta, uma interface web de manipulação foi criada para que o usuário possa interagir intuitivamente com o modelo sem precisar programar rotinas, bastando apenas informar com quais elementos deseja trabalhar.

Utilizando a linguagem de programação \texttt{Python} e o framework Flask, entre outras bibliotecas, além de \texttt{shell script}, foi possível receber do usuário o conjunto de dados, a identificação das colunas (variável resposta, tempo e indivíduo/tratamento) e a distribuição escolhida para o ajuste do modelo, retornando um arquivo \texttt{.zip} com todas as análises para cada restrição em função do tempo \textit{(Vide as figuras \ref{fig24cap23}, \ref{fig24cap24} e \ref{fig24cap25})}. 

% \begin{figure}[h]
% 	\centering
% 	\includegraphics[width=17cm]{tela1}
% 	\caption{Interface gráfica - Janela 01}
% 	\label{fig24cap23}
% \end{figure} 

% \begin{figure}[h]
% 	\centering
% 	\includegraphics[width=15cm]{tela2}
% 	\caption{Interface gráfica - Janela 02}
% 	\label{fig24cap24}
% \end{figure} 

% \begin{figure}[h]
% 	\centering
% 	\includegraphics[width=15cm]{tela3}
% 	\caption{Interface gráfica - Janela 03}
% 	\label{fig24cap25}
% \end{figure} 

\subsection{Resultados alcançados}
\label{cap3:secResultadosAlcancados}

Finalmente, chegou-se a um modelo de interação onde o pesquisador passou a ser um usuário dos modelos de regressão que descendem do GAMLLS. Com a retirada de ruído de informações que não dizem respeito a restrição da variável resposta em função do tempo, chegou-se a um esquema em tabela semelhante ao mostrado em \ref{tab8cap5histDist3}, onde:

\begin{itemize}
	\item "Época" refere-se a variável tempo no momento da comparação;
	\item "Rest." refere-se ao tratamento em restrição para o tempo $\omega_{i}$;
	\item "Comp." refere-se ao tratamento comparado com restrição para o tempo $\omega_{i}$;
	\item "Value", "Std Error", "DF", "t.value" e "p.value" são os valores retornados pelo modelo ajustado; e
	\item "Sign." é marcador da comparação que verifica se dentro do tempo $\omega_{i}$ a interação entre os dois tratamentos foi significativa ou não. \textbf{*} - foi significante a interação; \textbf{NS} - não foi significante a interação;
\end{itemize}


\small{
\begin{table}[h]
	\begin{center}
		\caption{Representação da saída de dados pela nova proposta}\vspace*{0.3cm}
		\label{tab8cap5histDist3}\vspace*{0.3cm}
		\begin{tabular}{lllllllll}
			Época & Rest. & Comp. & Value & Std Error & DF & t.value & p.value & Sign. \\
			\hline\hline
			$\omega_{0}$            & $y_{111}$            &  $y_{111}$      & -    & -       & -   & -        & -        & [ * or \texttt{NS} ]\\
			. \\
			. \\
			. \\
			$\omega_{i}$            & $y_{ijk}$            &  $y_{ijk}$      & -    & -       & -   & -        & -        & [ * or \texttt{NS} ]              
		\end{tabular}
	\end{center}
\end{table}
}

\section{Observações finais}
\label{cap4:observacoesFinais}

Neste capítulo foram analisados seis conjuntos de dados com variações diferentes para o comportamento da distribuição. O conjunto \textbf{LV-C}, assim como o \textbf{LV-COD}, apresentava assimetria a esquerda sendo o modelo OLLST aquele que saiu-se melhor nos testes comparativos em relação o NO. Há de se considerar, no entanto, que mesmo apresentando o melhor resultado entre os dois, tanto o OLLST quanto o NO, tiveram dificuldades para estimar os dados de \textbf{LV-COD}, de modo que ambos tiveram uma taxa de erro superior a 50\%. 

Quanto aos dados para o teor de carbono em razão do nitrogênio (conjunto \textbf{LV-CN}), considerando que o comportamento da distribuição assemelha-se ao gaussiano, o modelo NO obteve ligeira vantagem com todos os pontos dentro do intervalo de bandas do envelope simulado, ao passo que o modelo OLLST teve um apenas, dos 144 pontos fora desse intervalo. Sob essa perspectiva é válido afirmar que ambos os modelos saíram-se bem para o ajuste deste conjunto de dados.

Os conjunto de dados \textbf{LV-EN} e \textbf{LV-EC} apresentaram um comportamento semelhante em suas distribuições dos dados. A variação acentuada dos valores para a variável resposta pode ter contribuído para maior taxa de erro do OLLST em relação ao NO. Para os dados \textbf{LV-EN} a diferença, ainda que significativa, foi menor quando compara a do \textbf{LV-EC}. 

Por fim, na análise do conjunto \textbf{LV-N} o modelo OLLST obteve excelentes resultados, alcançando um ajuste significativamente melhor em relação ao NO. O gráfico de probabilidade mostrou uma taxa de efetividade de 92,36\% para o modelo, enquanto o NO manteve-se em 47,92\%. 

Neste capítulo, foram apresentados também os resultados práticos da melhoria proposta na seção \ref{cap3:secTratamentoSaidaGamlss}. A ferramenta desenvolvida com uso de diversas tecnologias pode representar um avanço na análise dos resultados obtidos pelos ajustes de modelos hierárquicos da biblioteca \texttt{gamlss}, ressaltando ainda que os dados apresentados na tabela \ref{tableResultados11} são frutos do ajuste proposto pela nova ferramenta.

Apesar das facilidades proposta pelo novo modelo de saída, existem ainda aspectos de melhorias que podem ser melhor explorados. Um exemplo disso seria a inclusão com extensão do editor de código para $\mathcal{R}$, \textit{rStudio}, agregação de ouros modelos ou bibliotecas, além do modelo hierárquico, entre outras. 


\chapter{Considerações Finais}
\label{cap5:title}

\section{Conclusões}

Como resultado das limitações observadas pelos autores \cite{azzalini1999statistical, braga2022random} para ajuste de modelos de regressão cujo os dados não apresentam distribuição gaussiana, somado aos desafios da análise de dados longitudinais \cite{fitzmaurice2008longitudinal}, conduziu-se uma pesquisa com o modelo OLLST a fim de analisa-lo sob essa perspectiva, até o momento não estudada.  

Duas foram as frentes de trabalho neste estudo: a primeira referia-se à análise do modelo OLLST, proposto por \cite{adolf2021}, para conjunto de dados com medida repetida no tempo; e a outra foi a proposta de melhoria da saída de dados para modelos pertencentes a família de distribuição \texttt{gamlss} \cite{rigby2005generalized}, incluindo assim o modelo OLLST.

Para análise do modelo com dados longitudinais foram considerados seis conjuntos de dados derivados estudo proposto por \cite{gabriel2022} sobre análise do teor de carbono no solo. A pesquisa comparou o desempenho do modelo OLLST com o NO para todos os conjuntos percebendo desempenho superior do OLLST para três, performance semelhante em um, e inferior em dois. 

Para os conjuntos \textbf{LV-CN}, \textbf{LV-EN} e \textbf{LV-EC}, o modelo OLLST apresentou uma performance menor que o NO, alguns caso, ligeiramente abaixo e em outros essa diferença foi um pouco mais acentuada. Nestas análises, foram verificadas a presença de picos dissonantes (\textit{outliers}) para variável reposta em relação aos demais valores, fator que pode ter influência na caracterização da amostra levando o modelo à estimativas menos acertivas. 

Já para o conjunto \textbf{LV-N}, com indicativo de bimodalidade, o OLLST ajustou-se significativamente melhor quando comparado com NO. Taxa um pouco mais modesta, ainda sim melhor que o NO, ocorreu na análise do conjunto \textbf{LV-C}, com tendência à assimetria a esquerda, onde apenas 36 dos 144 pontos do gráfico de probabilidade ficaram fora dos limites das banda, contra 44 do NO. 

Como resultado dos testes, é possível concluir pela validade do modelo OLLST para análise de dados longitudinais com efeito aleatório, como mostra a análise descritiva proposta no capitulo anterior para o conjunto \textbf{LV-C}, além do excelente ajuste do modelo para o conjunto \textbf{LV-N}. A presença de \textit{outliers}, trouxe um indicativo de que o modelo NO pode estimar melhor, em cenários assim, no entanto, esta deficiência não representa uma característica exclusiva do OLLST, sendo inclusive uma dificuldade já mapeada na literatura da análise de regressão para estimação (veja \cite{barnett1994outliers}). Na análise do dados \textbf{LV-COD}, por exemplo, mesmo com presença de \textit{outliers} o modelo OLLST conseguiu adequar-se melhor ao comportamento da distribuição quando comparado o NO, ainda que ambos não tenham conseguido encaixar mais de 50\% dos pontos dentro do intervalo do gráfico de probabilidade. 

Para o segundo nicho de trabalho desta pesquisa, como resultado, chegou-se a uma ferramenta cujo objetivo é tornar mais intuitivo o acesso a mecanismos de modelagem para dados com medida repetida no tempo. Com a rotina de ajuste do modelo construída, foi possível interagir com a linguagem $\mathcal{R}$ e obter resposta para modelagem dos dados sem fosse necessário codificar na linguagem, uma proposta que pode  difundir o acesso estatística.  

%Os testes primários com a nova distribuição Odd Log-Logistic Skew t-Student mostraram-se satisfatórios quanto na análise de dados longitudinais em modelos hierárquicos. O modelo ajustou-se bem ao conjunto de dados estudado e possui uma proposta de estimação de parâmetros robusta, se valendo método da máxima verossimilhança e o uso da quadratura de Gauss-Hermite para o cálculo das integrais.
%
%Tanto as análises de resíduos quanto os resultados obtidos pelo teste comparativo dois a dois, sugerem que o modelo possui capacidade suficiente para trabalhar com dados de medida repetida no tempo, não exigindo a distribuição normal dos dados, como a maioria dos modelos difundidos atualmente. Ainda assim, há pontos que necessitam de um olhar mais aprofundado, sobre tudo pela ótica de pesquisadores que não têm familiaridade com as bibliotecas $\mathcal{R}$ de análise de dados.

\section{Perspectiva de trabalhos futuros}
Ainda que a análise feita pelo OLLST tenha sido mais assertivas para alguns conjuntos de dados, existem lacunas que podem ser melhores exploradas para uma variedade maior de comportamento. Um exemplo disso, seria a condução de um estudo de simulação para o modelo tomando como base os parâmetros dispostos aqui e avaliar a eficiência das estimativas do modelo. 

Há também aperfeiçoamentos que podem atingir a ferramenta proposta neste trabalho. Por ser uma ferramenta web, ela pode ser propagada na internet ou incorporar como extensão do \texttt{rStudio}, podendo ainda agregar outras famílias de distribuição e bibliotecas, ou mesmo refinando ainda mais a saída para somente dados de interesse do pesquisador.  



% --- -----------------------------------------------------------------
% --- Referencias Bibliograficas. (Obrigatorio)
% --- -----------------------------------------------------------------
    \cleardoublepage
%    \bibliographystyle{acm-2} % abbrv - abnt-num
	%\bibliographystyle{plainnat}
    \bibliography{bibliografia} % arquivo fonte com a bibilografia

% --- -----------------------------------------------------------------
% --- Apendice.(Opcional)
% --- -----------------------------------------------------------------
\end{document}
